% !TeX root = ../../Book4.tex
% 纲要标题：## 第2章　极限、余极限与交换性质
% 纲要位置：计划纲要.md，第 3128 行。

\chapter{极限、余极限与交换性质}
\label{b4:ch:02}
\BookChapterNumber{b4:ch:02}

\begin{chapterabstract}
本章在集合、群、交换环与模中构造极限和余极限，证明由积与等化子构造小极限的存在判据，并比较滤过余极限、逆极限以及无限积、无限余积的正合性质。
\end{chapterabstract}

\begin{learninggoals}
\begin{itemize}
\item 根据泛性质区分积与余积、纤维积与纤维余积，并在不同代数范畴中计算它们。
\item 把一个相容映射问题写成图表及其锥，用积与等化子实际构造一般极限。
\item 写出模的正向极限和逆极限的元素描述，证明滤过余极限的正合性与逆极限的左正合性。
\item 写出函子保持极限的比较映射，说明滤过余极限与有限极限交换所需的两个条件。
\item 识别正合性对有限构造的作用，并解释无限积、无限余积中的反例。
\end{itemize}
\end{learninggoals}

给定两个集合映射 \(f:A\rightarrow C\) 与 \(g:B\rightarrow C\)。仅取笛卡儿积 \(A\times B\)，其中的数对未必满足 \(f(a)=g(b)\)；把不满足等式的数对剔除，所得集合便能把两个来源的数据按共同的像配对。另一方面，若要把 \(A\) 与 \(B\) 中指定的元素真正粘成一个对象，就不能只取不交并，还须施加关系。前一种构造由相容方程挑选元素，后一种构造由关系识别元素；本章将证明，它们分别体现极限与余极限的泛性质，并在模等具体范畴中追问这些构造怎样作用于正合列。

\ProseParagraphBreak
本章以\chapref{b4:ch:01}的范畴、函子和自然变换为基础。模的核、商、直和及短正合列沿用第二卷第3章的基本构造；交换环的纤维余积还用到第二卷第6.1节中张量积的平衡映射泛性质。这里只在具体模范畴中讨论正合性，一般交换范畴的公理安排在\chapref{b4:ch:04}。

\ProseParagraphBreak
模中的显式公式使极限与余极限的泛性质更容易辨认。正向极限中的一个元素可由某一阶段的元素表示，逆极限中的一个元素却是一整族相容数据；这一区别直接影响正合性。关于图表、极限与滤过范畴的定义及其相互关系，可查 Weibel 的《An Introduction to Homological Algebra》\cite[Section A.5, A.5.1--A.5.4, pp. 427--429]{Weibel1994HomologicalAlgebra}。

% !TeX root = ../../Book4.tex
% 纲要标题：### 2.1 基本极限与余极限的构造
% 纲要位置：计划纲要.md，第 2479 行。

\section{基本极限与余极限的构造}
\label{b4:sec:0201}

积和核都可以通过进入它们的映射来刻画；余积和商则通过从它们出发的映射来刻画。这种表达使同一个构造问题能够在群、环和模中分别求解。

% 纲要内容（仅作写作计划，尚非教材正文）：
%
% * 积、余积、等化子与余等化子
% * 纤维积与纤维余积
% * 在群、环、模中的不同计算方式
%

\subsection{积、余积、等化子与余等化子}
\label{b4:subsec:0201:01}

给定两个对象，常要同时记录通向它们的两个映射；给定两个平行映射，则常要找出使这两个映射相等的部分。积与等化子分别处理这两种要求。把全部箭头反向，又得到余积与余等化子。以下定义不预先规定对象必须由哪些元素组成，而是规定它应当容许哪些映射。

\begin{definition}\label{b4:def:initial-terminal-zero}
设 \(\mathcal C\) 为范畴。若对象 \(I\) 到每个对象 \(X\) 恰有一个态射，则称 \(I\) 为\term[shi-duixiang]{始对象}[initial object]，也称初对象。若每个对象 \(X\) 到 \(T\) 恰有一个态射，则称 \(T\) 为\term[zhong-duixiang]{终对象}[terminal object]。同时为始对象与终对象的对象称为\term[ling-duixiang]{零对象}[zero object]，常记为 \(0\)。
\end{definition}

集合范畴的始对象是空集，终对象是单点集，所以没有零对象。模范畴的零模则同时满足两种唯一性。任意两个始对象之间有唯一同构：分别取两个方向的唯一态射，其复合只能是各自的恒等态射；终对象同理。因此这些名称确定对象的泛性质，而不要求所选对象在字面上唯一。

\begin{definition}[积与余积]\label{b4:def:product-coproduct}
设 \(\mathcal C\) 为范畴，\((X_i)_{i\in I}\) 为集合 \(I\) 指标的对象族。对象 \(P\) 连同态射 \(p_i:P\rightarrow X_i\) 若满足：对每个对象 \(T\) 及每族态射 \(f_i:T\rightarrow X_i\)，存在唯一的 \(f:T\rightarrow P\)，使 \(p_if=f_i\)，则称 \((P,(p_i))\) 为该对象族的\term[ji]{积}[product]，记为 \(\prod_{i\in I}X_i\)。

\ProseParagraphBreak
对象 \(Q\) 连同态射 \(\iota_i:X_i\rightarrow Q\) 若满足：对每个对象 \(T\in\mathcal C\) 及每族态射 \(g_i:X_i\rightarrow T\)，存在唯一态射 \(g:Q\rightarrow T\)，使 \(g\iota_i=g_i\)，则称它为该对象族的\term[yuji]{余积}[coproduct]，记为 \(\coprod_{i\in I}X_i\)。
\end{definition}

\[
\begin{tikzcd}[column sep=large,row sep=large]
&P\arrow[d,"p_i"]\\
T\arrow[ur,dashed,"\exists!\,f"]\arrow[r,"f_i"']&X_i
\end{tikzcd}
\qquad
\begin{tikzcd}[column sep=large,row sep=large]
X_i\arrow[r,"\iota_i"]\arrow[dr,"g_i"']&Q\arrow[d,dashed,"\exists!\,g"]\\
&T.
\end{tikzcd}
\]

这里连同映射一起给出的数据才是积或余积。仅有一个同构于积的对象，却没有指定投影，就不足以写出具体的泛性质。空族的积要求从每个对象到它恰有一个态射，因此是终对象；空族的余积相应是始对象。集合范畴中，单点集为终对象，空集为始对象。

\begin{proposition}[积与余积的唯一性]\label{b4:prop:product-unique}
设 \(\mathcal C\) 为范畴，\((X_i)_{i\in I}\) 为集合 \(I\) 指标的 \(\mathcal C\) 中的对象族。该对象族的任意两个积之间，存在唯一的保持全部投影的同构；余积也满足对偶的结论。若 \((P,(p_i:P\rightarrow X_i))\) 为该对象族的积，则对任意对象 \(T\in\mathcal C\) 及态射 \(f,g:T\rightarrow P\)，条件 \(p_if=p_ig\) 对所有 \(i\in I\) 成立蕴含 \(f=g\)。
\end{proposition}
\begin{proof}
对两个积 \((P,p_i)\)、\((P',p_i')\)，泛性质分别给出 \(u:P\rightarrow P'\)、\(v:P'\rightarrow P\)，满足 \(p_i'u=p_i\)、\(p_iv=p_i'\)。于是 \(vu\) 与 \(1_P\) 具有相同的全部投影，由唯一性相等。同理 \(uv=1_{P'}\)。这也证明了同构的唯一性。最后一个断言正是泛性质中唯一性的另一种表述。反转箭头得到余积的证明。
\end{proof}

\begin{definition}[等化子与余等化子]\label{b4:def:equalizer-coequalizer}
设 \(f,g:X\rightarrow Y\) 是范畴 \(\mathcal C\) 中的两个平行态射。态射 \(e:E\rightarrow X\) 若满足 \(fe=ge\)，且每个满足 \(fh=gh\) 的 \(h:T\rightarrow X\) 均唯一分解为 \(h=e\widetilde h\)，则称它为 \(f,g\) 的\term[denghuazi]{等化子}[equalizer]。态射 \(q:Y\rightarrow Q\) 若满足 \(qf=qg\)，且每个满足 \(hf=hg\) 的 \(h:Y\rightarrow T\) 均唯一分解为 \(h=\widetilde hq\)，则称它为 \(f,g\) 的\term[yudenghuazi]{余等化子}[coequalizer]。
\end{definition}

\[
\begin{tikzcd}[column sep=large,row sep=large]
T\arrow[d,dashed,"\exists!\,\widetilde h"']\arrow[dr,"h"]&&\\
E\arrow[r,"e"']&X\arrow[r,shift left=0.7ex,"f"]
\arrow[r,shift right=0.7ex,"g"']&Y
\end{tikzcd}
\]
\[
\begin{tikzcd}[column sep=large,row sep=large]
X\arrow[r,shift left=0.7ex,"f"]\arrow[r,shift right=0.7ex,"g"']&
Y\arrow[r,"q"]\arrow[dr,"h"']&Q\arrow[d,dashed,"\exists!\,\widetilde h"]\\
&&T.
\end{tikzcd}
\]

等化子总是单态射。事实上，若 \(ea=eb\)，则 \(a,b\) 都是态射 \(ea\) 的分解，故相等。对偶地，余等化子总是满态射。这些结论只用了泛性质，不需要把单态射解释为单射。

\begin{example}[集合与模中的四种构造]\label{b4:exm:elementary-limits}
集合的积是笛卡儿积，余积是不交并；\(f,g:X\rightarrow Y\) 的等化子是集合 \(\{x\in X\mid f(x)=g(x)\}\) 的包含映射，余等化子则是 \(Y\) 按包含所有关系 \(f(x)\sim g(x)\) 的最小等价关系所取的商。关系本身未必具有传递性，所以这里必须取其生成的等价关系。

\ProseParagraphBreak
设 \(R\) 为含幺环，模均为含幺左 \(R\) 模。在 \(R\text{-}\mathbf{Mod}\) 中，
\[
\prod_{i\in I}M_i=\{\,(m_i)_{i\in I}\mid m_i\in M_i\,\},\qquad
\coprod_{i\in I}M_i=\bigoplus_{i\in I}M_i.
\]
直和只允许有限个非零分量。由直和给出的映射为
\(\sum_i\iota_i(m_i)\mapsto\sum_i f_i(m_i)\)，有限支集保证右边有意义。对于两个模同态 \(f,g:M\rightarrow N\)，等化子为 \(\ker(f-g)\) 的包含映射，余等化子为 \(N\rightarrow N/\operatorname{im}(f-g)\)。商模的泛性质恰好说明，令 \(f\) 与 \(g\) 在商后相等，就应当把每个 \(f(m)-g(m)\) 置为零。
\end{example}

有限个模的直和同时也是积，两个泛性质却仍有不同的方向。无限多个模时，有限支集条件将这两个对象分开；这个区别将在\secref{b4:sec:0204}产生实质影响。

\subsection{纤维积与纤维余积}
\label{b4:subsec:0201:02}

若两个映射 \(f:X\rightarrow Z\)、\(g:Y\rightarrow Z\) 给出同一对象的两种信息，把 \(X\) 与 \(Y\) 简单放在一起还不够，还必须要求它们在 \(Z\) 中的像相同。纤维积就是附加这一相容条件的积。

\begin{definition}[纤维积与纤维余积]\label{b4:def:pullback-pushout}
设 \(f:X\rightarrow Z\)、\(g:Y\rightarrow Z\) 是范畴 \(\mathcal C\) 中的态射。对象 \(P\) 连同 \(p:P\rightarrow X\)、\(q:P\rightarrow Y\) 若满足 \(fp=gq\)，且每对满足 \(fa=gb\) 的 \(a:T\rightarrow X\)、\(b:T\rightarrow Y\) 唯一分解为 \(a=ph,b=qh\)，则称它为该对态射的\term[xianweiji]{纤维积}[pullback]，记为 \(X\times_ZY\)。

\ProseParagraphBreak
设 \(u:A\rightarrow X\)、\(v:A\rightarrow Y\) 为 \(\mathcal C\) 中的态射。对象 \(Q\) 连同态射 \(i:X\rightarrow Q\)、\(j:Y\rightarrow Q\) 若满足 \(iu=jv\)，且对每个对象 \(T\in\mathcal C\) 及每对满足 \(au=bv\) 的态射 \(a:X\rightarrow T\)、\(b:Y\rightarrow T\)，存在唯一态射 \(h:Q\rightarrow T\)，使 \(a=hi,b=hj\)，则称它为这两个态射的\term[xianweiyuji]{纤维余积}[pushout]，记为 \(X\amalg_A Y\)。
\end{definition}

\[
\begin{tikzcd}[column sep=large,row sep=large]
T\arrow[r,dashed,"\exists!\,h"]\arrow[rr,bend left=22,"a"]
\arrow[dr,bend right=18,"b"']&P\arrow[r,"p"]\arrow[d,"q"']&X\arrow[d,"f"]\\
&Y\arrow[r,"g"']&Z
\end{tikzcd}
\]
\[
\begin{tikzcd}[column sep=large,row sep=large]
A\arrow[r,"u"]\arrow[d,"v"']&X\arrow[d,"i"]\arrow[dr,bend left=18,"a"]&\\
Y\arrow[r,"j"']\arrow[rr,bend right=22,"b"']&Q\arrow[r,dashed,"\exists!\,h"]&T.
\end{tikzcd}
\]

纤维余积回答的是另一种问题：在两个对象中已经指定了来自 \(A\) 的数据，怎样把它们粘在一起，使这些数据的两个像一致。定义并不要求 \(u,v\) 为单态射，也不保证 \(i,j\) 为单态射。

\begin{proposition}[模的纤维积与纤维余积]\label{b4:prop:module-pullback-pushout}
设 \(R\) 为含幺环，\(X,Y,Z,A\) 为左 \(R\) 模，所有映射均为 \(R\) 模同态。对于 \(f:X\rightarrow Z\)、\(g:Y\rightarrow Z\)，有
\[
X\times_ZY=\{\,(x,y)\in X\oplus Y\mid f(x)=g(y)\,\}.
\]
对于 \(u:A\rightarrow X\)、\(v:A\rightarrow Y\)，有
\[
X\amalg_A Y=(X\oplus Y)/
\operatorname{im}\bigl(A\xrightarrow{(u,-v)}X\oplus Y\bigr).
\]
这些公式分别配以坐标投影和直和的两个包含映射诱导的商映射。
\end{proposition}
\begin{proof}
第一式是子模。若 \(fa=gb\)，则 \(t\mapsto(a(t),b(t))\) 的像落在该子模中；其两个投影决定它，故分解存在且唯一。

第二式中 \(i(u(a))=j(v(a))\)，因为两者之差为 \((u(a),-v(a))\) 的类。若 \(a_0:X\rightarrow T\)、\(b_0:Y\rightarrow T\) 满足 \(a_0u=b_0v\)，则同态 \((x,y)\mapsto a_0(x)+b_0(y)\) 在被商去的子模上为零，从而唯一下降到商模。其在 \(X,Y\) 上的限制又唯一决定它。
\end{proof}

\begin{example}[由相同商粘合两个整数]\label{b4:exm:integer-pullback}
取两个约化映射 \(\mathbb Z\rightarrow\mathbb Z/n\mathbb Z\)，其中 \(n\ge1\)。纤维积为
\[
\{(a,b)\in\mathbb Z^2\mid a\equiv b\pmod n\}
=\mathbb Z(1,1)\oplus\mathbb Z(n,0).
\]
两个投影均满射，但这个子模通常不等于整个 \(\mathbb Z^2\)。条件记录的是两个整数的剩余类必须相同。
\end{example}

集合中的纤维积有同样的元素公式。集合的纤维余积则从 \(X\amalg Y\) 出发，把 \(u(a)\) 与 \(v(a)\) 识别，并取生成的等价关系。映射出这个商集，恰好等于分别给出 \(X,Y\) 上的映射且在 \(A\) 上一致，这就验证了泛性质。由此也可看出，同一个泛性质在不同范畴中往往需要不同的具体商构造。

\subsection{群、环与模中的极限计算}
\label{b4:subsec:0201:03}

群、环和模的许多极限都能在笛卡儿积中取满足相容条件的子对象。余极限则通常需要先加入足够多的元素，再施加关系。以下比较说明，只记住集合范畴中的公式并不足以完成代数计算。

\begin{proposition}[群中的基本构造]\label{b4:prop:group-limits-colimits}
在群范畴 \(\mathbf{Grp}\) 中，积为直积，两个同态 \(f,g:G\rightarrow H\) 的等化子为子群 \(\{x\in G\mid f(x)=g(x)\}\)。余等化子为
\[
H\longrightarrow H/N,\qquad
N=\NormalClosure{\,f(x)g(x)^{-1}\mid x\in G\,},
\]
其中正规闭包取在 \(H\) 中。两个群的余积为自由积；同态 \(u:A\rightarrow G\)、\(v:A\rightarrow H\) 的纤维余积为
\[
(G*H)/\NormalClosure{\,u(a)v(a)^{-1}\mid a\in A\,},
\]
式中 \(u(a),v(a)\) 表示它们在自由积中的像。
\end{proposition}
\begin{proof}
直积与等化子的公式由逐分量群运算直接验证。若 \(q:H\rightarrow K\) 满足 \(qf=qg\)，则其核包含每个 \(f(x)g(x)^{-1}\)，从而包含它们的正规闭包。这等价于 \(q\) 唯一经过所写商群分解，证明了余等化子公式。

自由积可以直接由第一卷第8章的自由群泛性质和群展示构造：在 \(G,H\) 的底集合不交并上取自由群，再商去分别表达 \(G,H\) 内部乘法和单位元的关系。所得群接收来自 \(G,H\) 的同态；任意一对 \(G\rightarrow K,H\rightarrow K\) 使这些关系成立，故唯一给出商群到 \(K\) 的同态。这就是余积的泛性质。最后再令 \(u(a)\) 与 \(v(a)\) 相等，恰好添加公式中的正规关系，得到纤维余积。
\end{proof}

这里的自由积不同于直积。例如 \(C_2*C_2\) 的两个生成元可以分别作用为实直线上的反射 \(t\mapsto-t\)、\(t\mapsto2-t\)。它们的复合是非平凡平移，具有无限阶。因此 \(C_2*C_2\) 是无限群，而 \(C_2\times C_2\) 只有四个元素。

\begin{proposition}[交换环中的基本构造]\label{b4:prop:ring-limits-colimits}
在允许零环、态射保持单位元的交换环范畴 \(\mathbf{CRing}\) 中，积为逐分量运算的直积，等化子为两个同态相等的子环。对于交换环同态 \(f,g:A\rightarrow B\)，余等化子为
\[
B\longrightarrow B/(\,f(a)-g(a)\mid a\in A\,).
\]
对于交换环 \(R\) 及两个交换 \(R\) 代数 \(A,B\)，相应的纤维余积为 \(A\otimes_RB\)。特别地，两个交换环的余积为 \(A\otimes_{\mathbb Z}B\)。始对象为 \(\mathbb Z\)，终对象为零环。
\end{proposition}
\begin{proof}
直积和等化子的单位元均来自相应的坐标单位元，所以所写构造仍在本范畴中。商环公式由一个同态满足 \(hf=hg\) 当且仅当其核包含全部差 \(f(a)-g(a)\) 得到。

在张量积上规定
\[
(a\otimes b)(a'\otimes b')=aa'\otimes bb',\qquad
1=1_A\otimes1_B.
\]
张量积的平衡关系保证乘法良定义；结合律与交换律可先对纯张量检查，再由双线性推广。给定两个环同态 \(\varphi:A\rightarrow T,\psi:B\rightarrow T\)，且它们在 \(R\) 上一致，则 \((a,b)\mapsto\varphi(a)\psi(b)\) 为 \(R\) 平衡的双加性映射，故唯一诱导加性映射 \(A\otimes_RB\rightarrow T\)。乘法公式说明它保持乘法与单位元。反之，任何兼容于两个结构映射的环同态必须把 \(a\otimes b=(a\otimes1)(1\otimes b)\) 送到这个乘积，因此唯一。最后，任何交换环唯一接收 \(\mathbb Z\) 的保幺同态，且唯一映到零环。
\end{proof}

张量积的模构造及平衡映射泛性质已在第二卷第6.1节建立。上面的证明另外核对了乘法，因此才能把模的构造用于交换环的纤维余积。对于一般非交换环，这个证明中的乘法公式与交换像条件不能照搬。

\begin{example}[余积可能成为零环]\label{b4:exm:ring-zero-coproduct}
在 \(\mathbf{CRing}\) 中，\(\mathbb Z/2\mathbb Z\) 与 \(\mathbb Z/3\mathbb Z\) 的余积为零环：任何同时接收它们的环中都有 \(2\cdot1=3\cdot1=0\)，于是 \(1=0\)。它们的积却是六元环 \(\mathbb Z/2\mathbb Z\times\mathbb Z/3\mathbb Z\)。若任意排除零环，这个例子就会失去所需的余积。
\end{example}

% !TeX root = ../../Book4.tex
% 纲要标题：### 2.2 一般极限
% 纲要位置：计划纲要.md，第 2485 行。

\section{一般极限}
\label{b4:sec:0202}

一族对象之间若有许多相容条件，逐条写等式容易漏掉所需箭头。用小范畴记录图表的形状，再用锥组织兼容映射，就能统一构造其极限。

% 纲要内容（仅作写作计划，尚非教材正文）：
%
% * 图表及其锥
% * 极限与余极限的泛性质
% * 存在性和由积、等化子构造的定理
%

\subsection{图表及其锥}
\label{b4:subsec:0202:01}

积要求给出一族映射，纤维积还要求两条复合相等。一般的极限只是把需要兼容的对象与箭头一起列出来。为防止“所有条件”成为含混的说法，先用一个小范畴记录图表的形状。

\begin{definition}[图表、锥与余锥]\label{b4:def:diagram-cone}
设 \(J\) 为小范畴，\(\mathcal C\) 为范畴。函子 \(D:J\rightarrow\mathcal C\) 称为一个 \(J\) 形\term[tubiao]{图表}[diagram]。给定对象 \(T\in\mathcal C\) 及一族态射 \(\lambda_j:T\rightarrow D(j)\)（\(j\in\operatorname{Ob}J\)）。若对每个 \(J\) 中的态射 \(\alpha:j\rightarrow k\) 都有
\[
D(\alpha)\lambda_j=\lambda_k.
\]
则称 \((T,(\lambda_j))\) 为 \(D\) 的一个\term[zhui]{锥}[cone]，称 \(T\) 为其顶点。对两个锥 \((T,(\lambda_j))\)、\((T',(\lambda_j'))\)，若态射 \(h:T\rightarrow T'\) 满足 \(\lambda_j'h=\lambda_j\) 对每个 \(j\) 成立，则称 \(h\) 为这两个锥之间的态射。

\ProseParagraphBreak
给定对象 \(T\in\mathcal C\) 及一族态射 \(\mu_j:D(j)\rightarrow T\)。若对每个 \(J\) 中的态射 \(\alpha:j\rightarrow k\) 都有 \(\mu_kD(\alpha)=\mu_j\)，则称 \((T,(\mu_j))\) 为 \(D\) 的一个\term[yuzhui]{余锥}[cocone]，称 \(T\) 为其顶点。对两个余锥 \((T,(\mu_j))\)、\((T',(\mu_j'))\)，若态射 \(h:T\rightarrow T'\) 满足 \(h\mu_j=\mu_j'\) 对每个 \(j\) 成立，则称 \(h\) 为这两个余锥之间的态射。
\end{definition}

\[
\begin{tikzcd}[column sep=2.5em,row sep=large]
&T\arrow[dl,"\lambda_j"']\arrow[dr,"\lambda_k"]&\\
D(j)\arrow[rr,"D(\alpha)"']&&D(k)
\end{tikzcd}
\]
\[
\begin{tikzcd}[column sep=2.5em,row sep=large]
D(j)\arrow[rr,"D(\alpha)"]\arrow[dr,"\mu_j"']&&D(k)\arrow[dl,"\mu_k"]\\
&T&
\end{tikzcd}
\]

固定 \(D\) 后，锥及其态射组成一个范畴；余锥也一样。锥中的“相容”不仅涉及图中画出的某几条箭头，而是涉及指标范畴中的每条态射。当 \(J\) 由若干箭头生成时，核对这些生成箭头即可，因为复合条件随后自动成立。

\ProseParagraphBreak
记 \(\Delta T:J\rightarrow\mathcal C\) 为把每个对象送到 \(T\)、每个态射送到 \(1_T\) 的常值图表。利用\secref{b4:sec:0102}的自然变换，锥就是自然变换 \(\Delta T\Rightarrow D\)，余锥就是 \(D\Rightarrow\Delta T\)。自然性方块的交换条件，正是定义中所写的相容等式。

\begin{example}[几种指标范畴]\label{b4:exm:diagram-shapes}
若 \(J\) 只有一族对象及各自的恒等态射，锥就是通向这些对象的一族任意态射，因而积属于这种形状。若 \(J\) 由两个平行箭头 \(0\rightrightarrows1\) 组成，则锥相容条件要求两个复合相等。若 \(J\) 的形状为 \(0\rightarrow2\leftarrow1\)，锥就是纤维积定义中的相容映射对。空范畴也被允许：此时每个对象恰有一个空锥。
\end{example}

\subsection{极限与余极限的泛性质}
\label{b4:subsec:0202:02}

\begin{definition}[极限与余极限]\label{b4:def:limit-colimit}
设 \(J\) 为小范畴，\(\mathcal C\) 为范畴，\(D:J\rightarrow\mathcal C\) 为图表。该图表的锥范畴中的终对象称为 \(D\) 的\term[jixian]{极限}[limit]，记为 \((\lim_JD,(\pi_j))\)。具体地，每个锥 \((T,(\lambda_j))\) 唯一来自一个态射 \(h:T\rightarrow\lim_JD\)，使 \(\pi_jh=\lambda_j\)。

\ProseParagraphBreak
该图表的余锥范畴中的始对象称为 \(D\) 的\term[yujixian]{余极限}[colimit]，记为 \((\operatorname{colim}_JD,(\iota_j))\)。每个余锥 \((T,(\mu_j))\) 唯一来自一个态射 \(h:\operatorname{colim}_JD\rightarrow T\)，使 \(h\iota_j=\mu_j\)。
\end{definition}

\[
\begin{tikzcd}[column sep=2.5em,row sep=large]
&T\arrow[dl,"\lambda_j"']\arrow[dr,"\lambda_k"]
\arrow[d,dashed,"\exists!\,h" description]&\\
D(j)\arrow[rr,bend right=24,"D(\alpha)"']&
\lim_JD\arrow[l,"\pi_j"]\arrow[r,"\pi_k"']&D(k)
\end{tikzcd}
\]
\[
\begin{tikzcd}[column sep=2.5em,row sep=large]
&T&\\
D(j)\arrow[r,"\iota_j"']\arrow[ur,"\mu_j"]
\arrow[rr,bend right=24,"D(\alpha)"']&
\operatorname{colim}_JD\arrow[u,dashed,"\exists!\,h" description]&
D(k)\arrow[l,"\iota_k"]\arrow[ul,"\mu_k"']
\end{tikzcd}
\]

\begin{proposition}[唯一性与诱导映射]\label{b4:prop:limit-uniqueness-functoriality}
设 \(J\) 为小范畴，\(\mathcal C\) 为范畴，\(D,D':J\rightarrow\mathcal C\) 为图表。一个图表的任意两个极限之间存在唯一的保持结构映射的同构；余极限也如此。若已选定 \(D,D'\) 的极限及其结构映射 \(\pi_j:\lim_JD\rightarrow D(j)\)、\(\pi_j':\lim_JD'\rightarrow D'(j)\)，则自然变换 \(\eta:D\Rightarrow D'\) 唯一诱导态射
\[
\lim_J\eta:\lim_JD\longrightarrow\lim_JD',
\qquad
\pi_j'\,(\lim_J\eta)=\eta_j\pi_j.
\]
这些诱导映射保持恒等变换与复合。余极限具有相同方向的诱导映射与相同的复合性质。
\[
\begin{tikzcd}[column sep=large,row sep=large]
\lim_JD\arrow[r,dashed,"\exists!\,\lim_J\eta"]\arrow[d,"\pi_j"']&
\lim_JD'\arrow[d,"\pi_j'"]\\
D(j)\arrow[r,"\eta_j"']&D'(j).
\end{tikzcd}
\]
\end{proposition}
\begin{proof}
唯一性证明与积的证明一致：用两个泛性质得到互相的映射，其两个复合都是保持结构映射的自映射，从而必须为恒等映射。

自然性给出
\(D'(\alpha)\eta_j\pi_j=\eta_kD(\alpha)\pi_j=\eta_k\pi_k\)，
所以 \((\eta_j\pi_j)\) 是 \(D'\) 上的锥，极限泛性质给出所求映射。恒等情形与复合情形均可逐个复合到 \(\pi_j'\) 后检查，再用唯一性得到。余极限的证明反转全部箭头即可。
\end{proof}

在所讨论的图表中选定极限后，就可以把取极限看作函子。这里的选择只是为了写出具体对象；任何两种选择之间已有唯一的相容同构。涉及真类范围的所有图表时，不在没有说明的情况下另加一次真类指标的选择。

\ProseParagraphBreak
在局部小范畴中，极限泛性质还可以写成集合间的双射
\[
\mathcal C(T,\lim_JD)
\ \cong\
\{\text{以 }T\text{ 为顶点的 }D\text{ 的锥}\}.
\]
这把一个通向极限的态射解释为一组相容态射。随 \(T\) 改变，这个双射与预复合相容；随图表改变，又与上一命题的诱导映射相容。\secref{b4:sec:0204}将由此证明 Hom 函子保持极限。

\begin{definition}[完备与余完备]\label{b4:def:complete-cocomplete-category}
若范畴 \(\mathcal C\) 中每个小图表都有极限，则称 \(\mathcal C\) 为\term[wanbeifanchou]{完备范畴}[complete category]；若每个小图表都有余极限，则称它为\term[yuwanbeifanchou]{余完备范畴}[cocomplete category]。若只要求指标范畴的对象和态射均有限，分别称为有限完备和有限余完备。
\end{definition}

这里的完备性属于范畴的泛性质，与赋范空间或局部环的完备性是不同概念。定义也不要求真类大小的图表存在极限。

\subsection{由积与等化子构造极限}
\label{b4:subsec:0202:03}

一般图表可能含有很多相容条件，但每个条件仍只是两个复合相等。因此，可以先用积收集全部候选数据，再用一个等化子同时施加全部条件。

\begin{theorem}[由积与等化子构造极限]\label{b4:thm:limits-products-equalizers}
设 \(J\) 为小范畴，\(\mathcal C\) 为范畴，\(D:J\rightarrow\mathcal C\) 为图表。假设下列两个积和所需的等化子存在：
\[
P=\prod_{j\in\operatorname{Ob}J}D(j),\qquad
Q=\prod_{\alpha:j\rightarrow k\ \text{in }J}D(k).
\]
设 \(p_j:P\rightarrow D(j)\) 为投影。定义 \(s,t:P\rightarrow Q\)，使其在 \(\alpha:j\rightarrow k\) 分量上分别为
\[
s_\alpha=D(\alpha)p_j,\qquad t_\alpha=p_k.
\]
若 \(e:E\rightarrow P\) 是 \(s,t\) 的等化子，则 \((E,(p_je))\) 是 \(D\) 的极限。对偶地，若相应余积与余等化子存在，则 \(D\) 的余极限是下列两个态射的余等化子：
\[
\coprod_{\alpha:j\rightarrow k}D(j)
\ \rightrightarrows\
\coprod_{j\in\operatorname{Ob}J}D(j),
\]
其中 \(\iota_j:D(j)\rightarrow\coprod_{j\in\operatorname{Ob}J}D(j)\) 为余积的结构映射；在 \(\alpha\) 对应的余积项上，两个态射分别为 \(\iota_kD(\alpha)\) 与 \(\iota_j\)。
\end{theorem}
\begin{proof}
由 \(se=te\)，对每个 \(\alpha:j\rightarrow k\) 都有
\(D(\alpha)p_je=p_ke\)，故 \((p_je)\) 是锥。给定任意锥 \((T,(\lambda_j))\)，积的泛性质唯一给出 \(h:T\rightarrow P\)，使 \(p_jh=\lambda_j\)。锥的相容性说明 \(sh=th\)，因为二者在 \(Q\) 的每个分量上相等。因此 \(h\) 唯一经过 \(e\) 分解。这一分解就是所求的 \(T\rightarrow E\)。其唯一性先由 \(P\) 的投影确定到 \(P\) 的映射，再由等化子的唯一性确定到 \(E\) 的映射。

对于余极限，令 \(q:\coprod_jD(j)\rightarrow C\) 为所列态射对的余等化子。等式 \(q\iota_kD(\alpha)=q\iota_j\) 说明 \((q\iota_j)\) 是余锥。任意余锥先唯一给出余积到其顶点的映射，相容性又保证这个映射使两个平行态射相等，故唯一经过 \(q\) 分解。
\end{proof}

\begin{corollary}[完备性的构造判据]\label{b4:cor:completeness-criterion}
具有所有小积及等化子的范畴是完备的；具有所有小余积及余等化子的范畴是余完备的。把“小”改为“有限”，结论仍成立。对任意含幺环 \(R\)，集合范畴 \(\mathbf{Set}\)、群范畴 \(\mathbf{Grp}\)、允许零环及保幺同态的交换环范畴 \(\mathbf{CRing}\)，以及幺性左 \(R\) 模范畴 \(R\text{-}\mathbf{Mod}\) 均具有所有小极限；其中 \(\mathbf{Set}\) 与 \(R\text{-}\mathbf{Mod}\) 也具有所有小余极限。
\end{corollary}
\begin{proof}
第一部分直接应用上一构造，有限指标范畴只需要有限个积项或余积项。所列范畴中的积及等化子已在\secref{b4:sec:0201}给出，集合与模的任意小余积、余等化子也已给出。
\end{proof}

群和交换环也余完备，但本章不需要再为任意多个环建立一套生成关系记号。群的任意小余积可按有限个群时同样的自由群构造得到；交换环的任意小余积可在所有元素符号生成的交换多项式环中加入各环的加法、乘法和单位关系，再用余等化子构造一般余极限。这些关系均由集合指标给出，故仍是通常的环构造。

\begin{example}[集合极限与余极限的统一公式]\label{b4:exm:set-limit-formulas}
对于小图表 \(D:J\rightarrow\mathbf{Set}\)，极限由相容元组组成：
\[
\lim_JD
=\{\,(x_j)\in\prod_jD(j)\mid
D(\alpha)(x_j)=x_k\text{ 对每个 }\alpha:j\rightarrow k\,\}.
\]
余极限则是 \(\coprod_jD(j)\) 对关系
\((j,x)\sim(k,D(\alpha)(x))\) 生成的等价关系所取的商。一般情况下，两个元素相等可能需要沿若干正反方向的关系反复连接，不能直接断言它们在某一个后续对象中已经相等。下一节将说明，滤过性恰好使这种更简单的判据成立。
\end{example}

% !TeX root = ../../Book4.tex
% 纲要标题：### 2.3 正向与逆向系统
% 纲要位置：计划纲要.md，第 3143 行。

\section{正向与逆向系统}
\label{b4:sec:0203}

递增子模的并允许先把有限个元素放入共同阶段再计算；逆系统则要求全部坐标同时相容。这一区别决定了两种极限的元素公式和后面的正合性质。

% 纲要内容（仅作写作计划，尚非教材正文）：
%
% * 滤过余极限即通常的正向极限
% * 逆极限及其元素描述；左正合性的相容提升证明
% * 模范畴中的显式公式
%

\subsection{滤过余极限与正向极限}
\label{b4:subsec:0203:01}

考虑递增的子模族。它们的并仍是子模，因为任意两个元素总能放入同一个较大的子模，再在那里完成加法。这个简单事实指出了一种重要条件：有限次计算所涉及的数据，应当能够在同一个对象中比较。

\begin{definition}[滤过范畴与正向系统]\label{b4:def:filtered-category}
设 \(J\) 为非空小范畴。若 \(J\) 满足以下两个条件：
\begin{enumerate}
\item 对任意对象 \(i,j\)，存在对象 \(k\) 及态射 \(i\rightarrow k,j\rightarrow k\)；
\item 对任意平行态射 \(u,v:i\rightarrow j\)，存在 \(w:j\rightarrow k\)，使 \(wu=wv\)。
\end{enumerate}
则称 \(J\) 为\term[lvguofanchou]{滤过范畴}[filtered category]。以滤过范畴为指标的余极限称为\term[lvguoyujixian]{滤过余极限}[filtered colimit]。设 \(\mathcal C\) 为范畴，\(I\) 为非空有向偏序集，即任意两个元素有共同上界。给定 \(\mathcal C\) 中的对象 \(M_i\) 及态射 \(f_{ij}:M_i\rightarrow M_j\)（\(i\le j\)）；若满足
\[
f_{ii}=1,\qquad f_{jk}f_{ij}=f_{ik}\quad(i\le j\le k)
\]
则称这些对象与态射组成一个\term[zhengxiangxitong]{正向系统}[direct system]。该系统的余极限也称为\term[zhengxiangjixian]{正向极限}[direct limit]，记为 \(\varinjlim_i M_i\)。
\end{definition}
\symidx{\(\varinjlim_i M_i\)：正向系统的余极限}

偏序范畴中两个给定对象之间至多有一条态射，所以滤过定义中的第二个条件自动成立。对于一般范畴，这个条件不能省略：它用来协调同一数据被两种不同路径送到同一对象时产生的差别。例如，由两个对象与两条不同的非恒等态射组成的范畴
\[
\begin{tikzcd}[column sep=large]
0\arrow[r,shift left=0.6ex,"u"]\arrow[r,shift right=0.6ex,"v"']&1
\end{tikzcd}
\]
满足共同目标条件，因为两个对象都能映到 \(1\)。但从 \(1\) 出发只有恒等态射，因而无法通过后复合使 \(u,v\) 相等；它不满足第二条条件，并不是滤过范畴。

\begin{lemma}[有限数据的共同比较]\label{b4:lem:filtered-finite-data}
设 \(J\) 为滤过范畴。从 \(J\) 中指定有限个对象以及它们之间有限条态射后，存在一个对象 \(k\) 及从这些对象到 \(k\) 的态射，使每条指定态射对应的三角形交换。
\end{lemma}
\begin{proof}
先反复使用共同目标条件，把所有指定对象映到一个对象 \(k_0\)，记已选的映射为 \(a_i:i\rightarrow k_0\)。对一条指定态射 \(\alpha:i\rightarrow j\)，要协调的是从 \(i\) 到 \(k_0\) 的两条态射 \(a_i\) 与 \(a_j\alpha\)：
\[
\begin{tikzcd}[column sep=large,row sep=large]
i\arrow[r,"\alpha"]\arrow[dr,"a_i"']&j\arrow[d,"a_j"]&\\
&k_0\arrow[r,"w"]&k_1.
\end{tikzcd}
\]
其中三角形起初未必交换。滤过的第二条件给出 \(w:k_0\rightarrow k_1\)，使 \(wa_i=wa_j\alpha\)。把所有已选映射都后复合 \(w\)，这一条件便成立，先前已成立的等式也仍然成立。对剩余的有限条态射重复此步即可。没有指定对象时，使用 \(J\) 非空即可。
\end{proof}

\begin{proposition}[滤过余极限的元素判据]\label{b4:prop:filtered-colimit-elements}
设 \(J\) 为滤过小范畴，\(D:J\rightarrow\mathbf{Set}\) 为图表。余极限的每个元素可由某个 \(x\in D(i)\) 表示。两个代表 \(x\in D(i)\)、\(y\in D(j)\) 给出同一元素，当且仅当存在 \(u:i\rightarrow k\)、\(v:j\rightarrow k\)，使 \(D(u)x=D(v)y\)。特别地，若 \(x,y\in D(i)\)，则它们在余极限中相等，当且仅当存在一个态射 \(w:i\rightarrow k\)，使 \(D(w)x=D(w)y\)。
\end{proposition}
\begin{proof}
在不交并上，把有这样一个共同目标的两个代表记为 \(x\sim y\)。反身性与对称性直接成立。为证传递性，设 \(x\in D(i)\)、\(y\in D(j)\)、\(z\in D(r)\)，并有
\[
\begin{gathered}
u:i\rightarrow k,\quad v:j\rightarrow k,\qquad D(u)x=D(v)y,\\
u':j\rightarrow l,\quad v':r\rightarrow l,\qquad D(u')y=D(v')z.
\end{gathered}
\]
先用第一个滤过条件选取 \(a:k\rightarrow t\)、\(b:l\rightarrow t\)。从 \(j\) 到 \(t\) 的两条路径 \(av\) 与 \(bu'\) 未必相等，故再用第二个条件选取 \(w:t\rightarrow t'\)，使 \(wav=wbu'\)。于是
\[
D(wau)x=D(wav)y=D(wbu')y=D(wbv')z.
\]
态射 \(wau:i\rightarrow t'\) 与 \(wbv':r\rightarrow t'\) 给出 \(x\sim z\) 的见证，因此 \(\sim\) 是等价关系。

这个关系包含每个图表关系 \(x\sim D(\alpha)x\)。反过来，若两个代表在共同目标中相等，则它们都由图表关系与那个共同像等价，所以 \(\sim\) 恰好是一般余极限公式中的生成等价关系。最后，当两个代表同属 \(D(i)\) 时，共同目标判据可能先给出两条不同的态射 \(u,v:i\rightarrow k\)。将它们后复合为同一条态射，即得到最后一个结论。
\end{proof}

这个判据有两个经常使用的结果：有限多个元素可以同时选在某个共同对象中；有限多个已经在余极限中成立的等式，可以在继续向后映射一次后同时成立。第二句话需要先取各等式的见证，再用\cref{b4:lem:filtered-finite-data}协调这些见证，不能把不同等式的目标对象当作自动相同。

\begin{example}[递增的并与分母]\label{b4:exm:direct-limit-localization}
若 \(M_i\) 是同一模 \(M\) 的有向递增子模族，则其正向极限为 \(\bigcup_iM_i\)：相容同态族在这个并上逐元定义，交叠处的一致性由共同上界保证。

\ProseParagraphBreak
对素数 \(p\)，整数模的系统
\[
\mathbb Z\xrightarrow{\times p}\mathbb Z
\xrightarrow{\times p}\mathbb Z\xrightarrow{\times p}\cdots
\]
以 \(n=0,1,\ldots\) 编号，其正向极限为 \(\mathbb Z[1/p]\)。第 \(n\) 项的 \(a\) 对应 \(a/p^n\)，该对应与过渡映射相容；两个分数相等，当且仅当它们在某个更大指标处的整数代表相等。为验证余极限的泛性质，给定交换群 \(M\) 及同态族 \(\phi_n:\mathbb Z\rightarrow M\)，满足 \(\phi_{n+1}(pa)=\phi_n(a)\)，定义
\[
\Phi:\mathbb Z[1/p]\longrightarrow M,\qquad
\Phi(a/p^n)=\phi_n(a).
\]
若 \(a/p^n=b/p^m\)，取 \(\ell\ge n,m\)，则 \(p^{\ell-n}a=p^{\ell-m}b\)。由同态族的相容性，
\[
\phi_n(a)=\phi_\ell(p^{\ell-n}a)
=\phi_\ell(p^{\ell-m}b)=\phi_m(b),
\]
故 \(\Phi\) 良定义。把两个分数写到同一分母后，由 \(\phi_\ell\) 的加法性可知 \(\Phi\) 为同态。每个分数都有上述形式，所以与全部结构映射相容的同态只能是 \(\Phi\)，唯一性也成立。
\end{example}

\subsection{逆极限及其元素描述}
\label{b4:subsec:0203:02}

正向系统把信息逐步送到后面；逆向系统则记录同一个对象在不同精度下的投影。逆极限要求所有精度的选择同时相容，因此其元素通常是一整族数据，而不是某一个阶段的一个元素。

\begin{definition}[逆向系统与逆极限]\label{b4:def:inverse-system}
设 \(I\) 为非空有向偏序集，\(\mathcal C\) 为范畴。对象 \(M_i\) 及态射 \(\rho_{ij}:M_j\rightarrow M_i\)（\(i\le j\)）若满足
\(\rho_{ii}=1\)、\(\rho_{ij}\rho_{jk}=\rho_{ik}\)，则称为一个\term[nixiangxitong]{逆向系统}[inverse system]，即图表 \(I^{\mathrm{op}}\rightarrow\mathcal C\)。该图表的极限称为\term[nijixian]{逆极限}[inverse limit]，记为 \(\varprojlim_i M_i\)。
\end{definition}
\symidx{\(\varprojlim_i M_i\)：逆向系统的极限}

在集合、群、交换环或模中，逆极限均可写为相容元组：
\[
\varprojlim_iM_i=
\{\,(m_i)\in\prod_iM_i\mid
\rho_{ij}(m_j)=m_i\;\text{对所有}\;i\le j\,\}.
\]
代数运算逐分量进行。投影映射给出一个锥，而任何其他锥只能把元素送到其全部分量组成的元组，因而上述公式满足泛性质。

\begin{example}[整数的相容剩余类]\label{b4:exm:adic-inverse-limit}
对素数 \(p\)，约化映射组成逆向系统
\[
\cdots\longrightarrow\mathbb Z/p^3\mathbb Z
\longrightarrow\mathbb Z/p^2\mathbb Z
\longrightarrow\mathbb Z/p\mathbb Z.
\]
其逆极限环记为 \(\mathbb Z_p\)。一个元素是一族相容剩余类 \((a_n)\)，并不要求存在一个整数同时代表全部 \(a_n\)。每个整数给出这样的族；若其全部剩余类均为零，则它被所有 \(p^n\) 整除，故为零。因此 \(\mathbb Z\rightarrow\mathbb Z_p\) 为单射。

\ProseParagraphBreak
逐步选择代表可将每族唯一写成数字 \(c_0,c_1,\ldots\)，其中 \(0\le c_i<p\)，且
\[
a_n\equiv c_0+c_1p+\cdots+c_{n-1}p^{n-1}\pmod{p^n}.
\]
先由 \(a_1\) 决定 \(c_0\)，再由相容性逐项决定下一位即可。这里的无穷数字首先表达相容剩余类，不涉及实数意义下的无穷级数。
\end{example}
\symidx{\(\mathbb Z_p\)：\(p\)-进整数环}

\begin{example}[逆极限不保持满射]\label{b4:exm:inverse-limit-not-exact}
对 \(n\ge1\)，考虑短正合列
\[
0\longrightarrow2^n\mathbb Z\longrightarrow\mathbb Z
\longrightarrow\mathbb Z/2^n\mathbb Z\longrightarrow0.
\]
向较小指标的过渡映射依次为包含映射、恒等映射和约化映射，故各列组成相容的逆向系统。取逆极限后得到
\[
0\longrightarrow0\longrightarrow\mathbb Z
\longrightarrow\mathbb Z_2.
\]
给定 \((b_n)\in\mathbb Z_2\)，每个 \(b_n\) 单独都能选择整数代表 \(c_n\)。但中间系统的过渡映射是恒等映射，相容提升必须满足 \(c_1=c_2=\cdots\)，也就是由同一个整数代表全部 \(b_n\)。这样的整数未必存在，所以最后的映射不是满射。事实上，令 \(b_n\) 为 \(3\) 在
\(\mathbb Z/2^n\mathbb Z\) 中的逆元。逆元的唯一性说明 \((b_n)\) 相容。若它来自整数 \(b\)，则 \(3b-1\) 被所有 \(2^n\) 整除，故 \(3b=1\)，矛盾。
\end{example}

每个阶段的满射都能逐个提升元素，却未必能够把这些提升选成彼此相容的一族。上例显示，讨论逆极限的正合性时，不能把正向极限的结论直接反转。核中的元素则不同：它在对应的子模中有唯一前像，这种唯一性会迫使各阶段的前像相容。

\begin{proposition}[逆极限的左正合性]\label{b4:prop:inverse-limits-left-exact}
设 \(R\) 为含幺环，\(I\) 为非空小有向偏序集，\(A,B,C\) 为以 \(I\) 为指标的幺性左 \(R\) 模逆向系统。记它们的过渡同态分别为 \(\rho^A_{ij},\rho^B_{ij},\rho^C_{ij}\)（\(i\le j\)）。若同态族 \(u_i:A_i\rightarrow B_i\)、\(v_i:B_i\rightarrow C_i\) 与过渡同态相容，且各阶段的序列
\[
0\longrightarrow A_i\xrightarrow{u_i}B_i
\xrightarrow{v_i}C_i
\]
正合，则诱导的序列
\[
0\longrightarrow\varprojlim_iA_i
\xrightarrow{\overline u}\varprojlim_iB_i
\xrightarrow{\overline v}\varprojlim_iC_i
\]
也正合。因此逆极限函子是左正合的。
\end{proposition}
\begin{proof}
诱导同态逐分量施加 \(u_i\) 与 \(v_i\)。若相容族 \((a_i)\) 满足 \(u_i(a_i)=0\) 对所有 \(i\) 成立，则各 \(u_i\) 的单射性给出 \(a_i=0\)，所以 \(\overline u\) 为单射。由 \(v_iu_i=0\)，还有 \(\operatorname{im}\overline u\subseteq\ker\overline v\)。

反过来，若相容族 \((b_i)\in\varprojlim_iB_i\) 属于 \(\ker\overline v\)，则 \(v_i(b_i)=0\)。由 \(\operatorname{im}u_i=\ker v_i\)，存在 \(a_i\in A_i\) 使 \(u_i(a_i)=b_i\)；再由 \(u_i\) 单射，这个原像唯一。对 \(i\le j\)，过渡同态的相容性给出
\[
u_i\bigl(\rho^A_{ij}(a_j)\bigr)
=\rho^B_{ij}\bigl(u_j(a_j)\bigr)
=\rho^B_{ij}(b_j)
=b_i=u_i(a_i).
\]
由 \(u_i\) 单射，\(\rho^A_{ij}(a_j)=a_i\)。所以 \((a_i)\) 是 \(A\) 的相容族，且 \(\overline u((a_i))=(b_i)\)，得到中间位置的正合性。
\end{proof}

\subsection{模范畴中的显式公式}
\label{b4:subsec:0203:03}

模的极限通过方程截取直积，余极限则通过关系商去直和。这正是\cref{b4:thm:limits-products-equalizers}的加性形式：两个同态 \(s,t\) 的等化子是 \(\ker(s-t)\)，余等化子是 \(\operatorname{coker}(s-t)\)。下面的 \(\delta\) 是极限构造中两个直积同态之差，\(d\) 则是余极限构造中两个直和同态之差，把应当相等的两项之差置为零。这两种公式都适用于任意小指标范畴，不要求指标有向。

\begin{proposition}[模的极限与余极限公式]\label{b4:prop:module-limit-formulas}
设 \(R\) 为含幺环，\(J\) 为小范畴，\(M:J\rightarrow R\text{-}\mathbf{Mod}\) 为幺性左 \(R\) 模图表，记 \(M_j=M(j)\)。以下 \(\iota_j\) 与 \(\iota_\alpha\) 分别表示以对象和态射为指标的直和的结构包含。定义
\[
\begin{aligned}
\delta:\prod_{j\in\operatorname{Ob}J}M_j
&\longrightarrow\prod_{\alpha:j\rightarrow k}M_k,
&
\delta((m_j))_\alpha&=M(\alpha)m_j-m_k,\\
d:\bigoplus_{\alpha:j\rightarrow k}M_j
&\longrightarrow\bigoplus_{j\in\operatorname{Ob}J}M_j,
&
d(\iota_\alpha(m))&=\iota_k(M(\alpha)m)-\iota_j(m).
\end{aligned}
\]
则 \(\lim_JM=\ker\delta\)，而
\[
\operatorname{colim}_JM
=\operatorname{coker}d
=\left(\bigoplus_jM_j\right)/S,
\]
其中关系子模为
\[
S=\bigl\langle\,\iota_k(M(\alpha)m)-\iota_j(m)\mid \alpha:j\rightarrow k\;\text{在}\;J\;\text{中},\ m\in M_j\,\bigr\rangle_R.
\]
\end{proposition}
\begin{proof}
第一式的核正好由相容元组组成。第二式中，直和的每个元素只有有限支集，所以 \(d\) 良定义；其像是公式所列关系生成的子模。给定一族同态 \(f_j:M_j\rightarrow N\)，它们唯一给出直和上的同态。该同态在 \(\operatorname{im}d\) 上为零，当且仅当 \(f_kM(\alpha)=f_j\) 对所有 \(\alpha\) 成立。这正是余锥条件，因此它唯一下降到所列余核，得到泛性质。
\end{proof}

\begin{proposition}[滤过模图表的单阶段表示]\label{b4:prop:filtered-module-elements}
设 \(R\) 为含幺环，\(J\) 为滤过小范畴，\(M:J\rightarrow R\text{-}\mathbf{Mod}\) 为幺性左 \(R\) 模图表，记 \(M_i=M(i)\)。其模余极限的每个元素均可表示为某个 \(m\in M_i\) 的像。两个代表 \(m\in M_i\)、\(n\in M_j\) 给出同一元素，当且仅当存在 \(u:i\rightarrow k\)、\(v:j\rightarrow k\)，使 \(M(u)m=M(v)n\)。特别地，\(m\in M_i\) 的像为零，当且仅当某个 \(u:i\rightarrow j\) 满足 \(M(u)m=0\)。
\end{proposition}
\begin{proof}
先在底集合的滤过余极限上定义运算。对 \(m\in M_i,n\in M_j\)，选 \(u:i\rightarrow k\)、\(v:j\rightarrow k\)，定义
\[
[m]+[n]=[M(u)m+M(v)n].
\]
若另选 \(u':i\rightarrow l\)、\(v':j\rightarrow l\)，则\cref{b4:lem:filtered-finite-data}给出 \(a:k\rightarrow t\)、\(b:l\rightarrow t\)，使 \(au=bu'\)、\(av=bv'\)。同态的线性性给出
\[
M(a)\bigl(M(u)m+M(v)n\bigr)
=M(b)\bigl(M(u')m+M(v')n\bigr),
\]
所以两个共同目标算出的和表示同一个类。

若更换代表为 \(m'\in M_{i'}\)、\(n'\in M_{j'}\)，且 \([m]=[m']\)、\([n]=[n']\)，元素判据分别给出
\[
\begin{gathered}
p:i\rightarrow k_0,\quad p':i'\rightarrow k_0,\qquad M(p)m=M(p')m',\\
q:j\rightarrow k_1,\quad q':j'\rightarrow k_1,\qquad M(q)n=M(q')n'.
\end{gathered}
\]
选取 \(a:k_0\rightarrow t\)、\(b:k_1\rightarrow t\)，就有
\[
M(ap)m+M(bq)n=M(ap')m'+M(bq')n'.
\]
由已证明的共同目标无关性，可用 \(t\) 计算两组代表的和，因此更换代表也不改变结果。标量乘法定义为 \(r[m]=[rm]\)；代表在共同目标中的相等，经线性映射仍保持乘以 \(r\) 后的相等，故它也良定义。全部模公理只涉及有限多个元素，可送到一个共同对象后由该模中的公理验证。

所得模接收各 \(M_i\) 的相容线性映射；任意模值余锥在底集合层面诱导的唯一映射保持加法与标量乘法，所以也是模同态。因此这个模就是模范畴的余极限，底集合的元素判据随之成立。零元可在任意对象中取零代表，对同阶段的两个元素使用单个后续态射的判据，即得最后一句。
\end{proof}

\begin{proposition}[滤过余极限的正合性]\label{b4:prop:filtered-colimits-exact}
设 \(R\) 为含幺环，\(J\) 为滤过小范畴，\(A,B,C:J\rightarrow R\text{-}\mathbf{Mod}\) 为幺性左 \(R\) 模图表，记 \(A_j=A(j)\)、\(B_j=B(j)\)、\(C_j=C(j)\)。若自然变换 \(A\xrightarrow{u}B\xrightarrow{v}C\) 在每个对象 \(j\) 上给出短正合列
\[
0\longrightarrow A_j\xrightarrow{u_j}B_j
\xrightarrow{v_j}C_j\longrightarrow0,
\]
则诱导的序列
\[
0\longrightarrow\operatorname{colim}_JA
\xrightarrow{\overline u}\operatorname{colim}_JB
\xrightarrow{\overline v}\operatorname{colim}_JC
\longrightarrow0
\]
也是短正合列。
\end{proposition}
\begin{proof}
若 \(a\in A_i\) 的类被 \(\overline u\) 送到零，则存在 \(\alpha:i\rightarrow j\)，使
\(B(\alpha)u_i(a)=0\)。由自然性，这等于
\(u_jA(\alpha)(a)\)。\(u_j\) 为单射，故 \(A(\alpha)(a)=0\)，于是 \(a\) 在余极限中为零，证明了单射性。

若 \(b\in B_i\) 的类属于 \(\ker\overline v\)，则某个 \(\alpha:i\rightarrow j\) 使
\(v_jB(\alpha)(b)=0\)。阶段 \(j\) 的正合性给出 \(a\in A_j\)，满足 \(u_j(a)=B(\alpha)(b)\)，所以 \([b]=\overline u([a])\)。反向包含由 \(v_ju_j=0\) 得到。最后，任意 \(c\in C_i\) 可在同一阶段提升为某个 \(b\in B_i\)，故 \(\overline v\) 为满射。
\end{proof}

这项论证适用于任意过渡同态。关键在于：一个元素在余极限中为零，就能在某个后续对象中变为零。有限多个这样的等式还可放到同一阶段处理，这也将解释滤过余极限为什么能与有限极限交换。

% !TeX root = ../../Book4.tex
% 纲要标题：### 2.4 函子与极限
% 纲要位置：计划纲要.md，第 3149 行。

\section{函子与极限}
\label{b4:sec:0204}

把极限锥施加函子以后，得到的仍是一个锥，却未必仍有极限的泛性质。如果两边的极限都存在，就可以借助比较映射判断这一性质是否保持。有限构造与无限构造所需的保持条件不同，正合性本身只控制其中一部分。

% 纲要内容（仅作写作计划，尚非教材正文）：
%
% * 函子保持极限与同类极限的交换
% * 滤过余极限与有限极限交换；共同阶段中的代表元与相容等式
% * 模范畴中的左正合、右正合与正合函子；有限构造的保持性和无限构造的反例
%
% ---
%

\subsection{函子保持极限与同类极限的交换}
\label{b4:subsec:0204:01}

把一个图表先取极限再施加函子，与先施加函子再取极限，是两种不同的操作。即使两边都有对象，也还需要一个由泛性质决定的比较映射，才能准确询问它们是否相同。

\begin{definition}[保持极限与余极限]\label{b4:def:preserves-limits}
设 \(\mathcal C,\mathcal D\) 为范畴，\(F:\mathcal C\rightarrow\mathcal D\) 为函子，\(J\) 为小范畴，\(D:J\rightarrow\mathcal C\) 为图表，并选定极限锥 \((\lim_JD,(\pi_j:\lim_JD\rightarrow D(j)))\)。若 \(F\) 把此锥送到 \(FD\) 的极限锥，则称 \(F\) 保持这个极限。当 \(\lim_JFD\) 也已选定时，这等价于由锥 \(F(\pi_j)\) 诱导的比较映射
\[
F(\lim_JD)\longrightarrow\lim_JFD
\]
为同构。保持余极限的定义相应地要求余极限余锥被送到余极限余锥；在两边均存在时，比较映射的方向为
\[
\operatorname{colim}_JFD\longrightarrow
F(\operatorname{colim}_JD).
\]
\end{definition}

\[
\begin{tikzcd}[column sep=large,row sep=large]
F(\lim_JD)\arrow[r,dashed,"\exists!"]\arrow[dr,"F(\pi_j)"']&
\lim_JFD\arrow[d,"\text{投影}"]\\
&F(D(j))
\end{tikzcd}
\]
\[
\begin{tikzcd}[column sep=large,row sep=large]
F(D(j))\arrow[r,"\text{结构映射}"]\arrow[dr,"F(\iota_j)"']&
\operatorname{colim}_JFD\arrow[d,dashed,"\exists!"]\\
&F(\operatorname{colim}_JD).
\end{tikzcd}
\]

若 \(F\) 保持相关大小的积与等化子，则由\cref{b4:thm:limits-products-equalizers}可知它保持相应的极限：把该定理构造中的两个积及其等化子依次施加 \(F\)，所得仍是同一图表的极限构造。对偶地，保持余积与余等化子足以保证保持相应的余极限。这里包括空积或空余积，不能只检查非空族。

\begin{proposition}[Hom 函子保持极限]\label{b4:prop:representables-preserve-limits}
设 \(\mathcal C\) 为局部小范畴，\(X\in\mathcal C\)。函子
\(\operatorname{Hom}_{\mathcal C}(X,-):\mathcal C\rightarrow\mathbf{Set}\)
保持所有已经存在的小极限。具体地，设 \(J\) 为小范畴，\(D:J\rightarrow\mathcal C\) 为图表；若 \(D\) 有极限，则
\[
\operatorname{Hom}_{\mathcal C}(X,\lim_JD)
\ \cong\
\lim_{j\in J}\operatorname{Hom}_{\mathcal C}(X,D(j)).
\]
对偶地，若 \(D\) 有余极限，则
\[
\operatorname{Hom}_{\mathcal C}(\operatorname{colim}_JD,X)
\ \cong\
\lim_{j\in J^{\mathrm{op}}}\operatorname{Hom}_{\mathcal C}(D(j),X).
\]
\end{proposition}
\begin{proof}
第一式右边的元素是一族态射 \(f_j:X\rightarrow D(j)\)，满足
\(D(\alpha)f_j=f_k\)。它正是以 \(X\) 为顶点的锥，所以极限泛性质将它与唯一的 \(f:X\rightarrow\lim_JD\) 对应。这个双射就是比较映射，因为它把 \(f\) 送到 \((\pi_jf)\)。第二式的右边是与图表相容的态射族 \(D(j)\rightarrow X\)，余极限泛性质给出所求双射。两式都与预复合、后复合相容，因此也是自然的对应。
\end{proof}

这个证明只使用泛性质。\chapref{b4:ch:03}把这类 Hom 函子称为可表函子，并进一步研究哪些函子来自一个对象。

\begin{proposition}[同类极限的交换]\label{b4:prop:limits-fubini}
设 \(I,J\) 为小范畴，\(\mathcal C\) 为范畴，\(D:I\times J\rightarrow\mathcal C\) 为图表。若 \(\mathcal C\) 具有所需的小极限，则存在保持全部结构映射的自然同构
\[
\lim_{i\in I}\lim_{j\in J}D(i,j)
\ \cong\
\lim_{(i,j)\in I\times J}D(i,j)
\ \cong\
\lim_{j\in J}\lim_{i\in I}D(i,j).
\]
若 \(\mathcal C\) 具有所需的小余极限，则余极限之间也有相应的交换同构。
\end{proposition}
\begin{proof}
从对象 \(T\) 到第一个迭代极限的态射，等价于一族对 \(i\) 相容的态射
\(T\rightarrow\lim_jD(i,j)\)，又等价于一族
\(T\rightarrow D(i,j)\)，对 \(i,j\) 两个方向的箭头均相容。积范畴中的每条箭头都可分解为这两个方向的箭头，所以它正是积范畴图表上的锥。这就验证了该迭代对象满足中间极限的泛性质。交换 \(i,j\) 得到另一同构。所有同构由结构映射确定，因而具有所述自然性。反转箭头得到余极限版本。
\end{proof}

这一命题讨论极限与极限、余极限与余极限的交换。极限与余极限混合时，通常需要额外条件，下面的滤过性与有限性就是一组重要的充分条件。

\subsection{有限与无限的积和余积}
\label{b4:subsec:0204:02}

在有限积中选取元素，只涉及有限个分量；在等化子中检验元素，只涉及一个等式。因此，滤过余极限的单阶段判据可以处理有限极限。但对于无限多个分量，各分量需要的阶段可能没有共同上界。

\begin{proposition}[滤过余极限与有限极限交换]\label{b4:prop:filtered-colimits-finite-limits}
设 \(I\) 为滤过小范畴，\(J\) 为对象与态射均有限的范畴，\(\mathcal C\) 为 \(\mathbf{Set}\) 或 \(R\text{-}\mathbf{Mod}\)，其中 \(R\) 为含幺环。对每个图表 \(D:I\times J\rightarrow\mathcal C\)，自然比较映射
\[
\operatorname{colim}_{i\in I}\lim_{j\in J}D(i,j)
\longrightarrow
\lim_{j\in J}\operatorname{colim}_{i\in I}D(i,j)
\]
是同构。
\end{proposition}
\begin{proof}
先取 \(\mathcal C=\mathbf{Set}\)，且 \(J\) 非空。右边的一个元素是一族相容的类 \((\overline x_j)\)。为每个 \(\overline x_j\) 选取一个代表，设这些代表所在的对象为 \(i_1,\ldots,i_r\)。因为 \(J\) 只有有限个对象，滤过性给出一个共同目标 \(i\)，使全部 \(\overline x_j\) 均由 \(x_j\in D(i,j)\) 表示。接下来还要使这些代表满足全部相容等式；两个步骤所用的对象如下：
\[
\begin{tikzcd}[column sep=large,row sep=small]
i_1\arrow[dr]&\cdots&i_r\arrow[dl]\\
&i\arrow[d]&\\
&i'&
\end{tikzcd}
\]
先把有限多个代表送到 \(i\)，再把有限多个等式安排在 \(i'\) 中。具体地，对每条 \(\alpha:j\rightarrow k\)，相容性说明 \(D(i,\alpha)x_j\) 与 \(x_k\) 在滤过余极限中相等，因此存在一个从 \(i\) 出发的态射，使这两个元素的像相等。只有有限条箭头，由\cref{b4:lem:filtered-finite-data}可把这些见证协调成一个态射 \(i\rightarrow i'\)，使所有等式同时成立。于是这些像构成 \(\lim_jD(i',j)\) 的元素，证明比较映射满射。

若左边两个元素具有相同的像，先把它们的代表送到一个共同对象 \(i\)。这给出两个相容元组 \((x_j)\)、\((y_j)\)。右边相等表示每个 \(x_j,y_j\) 在滤过余极限中相等；有限多个等式可在同一个后续对象实现，因此两个元组在那里相等。它们在左边的类也相等，证明单射。

当 \(J\) 为空时，每个内层极限是单点集。非空滤过范畴中任意两个对象有共同目标，所以常值单点集图表的余极限仍是单点集，两边相同。

在模范畴中，极限的底集合由相容元组组成，滤过余极限的底集合由\cref{b4:prop:filtered-module-elements}给出同样的元素模型。因而刚才的论证证明比较映射为双射。它本来就是模同态，所以是模同构；空图表时两边均为零模。
\end{proof}

\begin{example}[无限积使比较映射失去满射性]\label{b4:exm:filtered-infinite-product}
设 \(k\) 为域，以非负整数 \(n\) 组成有向指标集，以正整数 \(j\) 标记一族模。令
\[
D_{n,j}=
\begin{cases}
k,&j\le n,\\
0,&j>n,
\end{cases}
\]
并以自然包含作为 \(n\) 方向的映射。对每个固定的 \(j\)，都有 \(\varinjlim_nD_{n,j}=k\)，但比较映射为
\[
\varinjlim_n\prod_{j\ge1}D_{n,j}
=\bigoplus_{j\ge1}k
\ \longrightarrow\
\prod_{j\ge1}\varinjlim_nD_{n,j}
=\prod_{j\ge1}k.
\]
它遗漏了 \((1,1,\ldots)\)。左边任何元素都在一个有限阶段出现，因此只有有限多个非零分量；右边允许各分量在不同阶段出现，没有统一的有限阶段。
\end{example}

有限个模的积与余积相同，不表示相应的无限构造也相同。一般的自然同态
\(\bigoplus_iM_i\rightarrow\prod_iM_i\)
是把有限支集元组视为任意元组；只有当其余元组都满足有限支集条件时，它才是满射。在集合范畴中，甚至两个非空集合的积与不交并也没有这样统一的同构。

\subsection{模范畴中的正合函子}
\label{b4:subsec:0204:03}

第二卷在模范畴中研究了正合性。现在可以问：保持短正合列，是否足以保持本章的全部构造？对于有限极限与有限余极限，答案是肯定的；对于无限构造，还需要更多条件。

\begin{definition}[模范畴之间的正合函子]\label{b4:def:exact-module-functor}
设 \(R,S\) 为含幺环，\(R\text{-}\mathbf{Mod}\) 与 \(S\text{-}\mathbf{Mod}\) 分别为幺性左 \(R\) 模与幺性左 \(S\) 模的范畴。给定函子 \(F:R\text{-}\mathbf{Mod}\rightarrow S\text{-}\mathbf{Mod}\)。若 \(F\) 在每个 Hom 群上保持加法，则称 \(F\) 为加性函子。

设 \(F\) 为加性函子。若对每个短正合列 \(0\rightarrow A\xrightarrow{u}B\xrightarrow{v}C\rightarrow0\)，所得序列
\[
0\longrightarrow F(A)\xrightarrow{F(u)}F(B)
\xrightarrow{F(v)}F(C)
\]
在 \(F(A)\) 与 \(F(B)\) 处正合，则称 \(F\) 为左正合函子；若所得序列
\[
F(A)\xrightarrow{F(u)}F(B)\xrightarrow{F(v)}F(C)
\longrightarrow0
\]
在 \(F(B)\) 与 \(F(C)\) 处正合，则称 \(F\) 为右正合函子。若 \(F\) 同时左正合与右正合，即把每个短正合列送到短正合列，则称 \(F\) 为正合函子。
\end{definition}

\begin{proposition}[正合性与有限构造]\label{b4:prop:exact-preserves-finite}
设 \(R,S\) 为含幺环。正合函子
\(F:R\text{-}\mathbf{Mod}\rightarrow S\text{-}\mathbf{Mod}\)
保持全部有限极限与有限余极限。
\end{proposition}
\begin{proof}
加性给出 \(F(0)=0\)：零模的恒等映射是零态射，故 \(F(0)\) 的恒等映射也为零，它只能是零模。对有限直和的包含与投影，关系
\(p_i\iota_j=\delta_{ij}\)、\(\sum_i\iota_ip_i=1\)
在施加加性函子后仍成立。这些关系表明 \(F\) 保持有限直和，也就是模中的有限积与余积。

对同态 \(f:M\rightarrow N\)，令 \(I=\operatorname{im}f\)，将它分解为满射 \(q:M\rightarrow I\) 与包含 \(j:I\rightarrow N\)，故 \(f=jq\)。分别对两个短正合列
\[
0\longrightarrow\ker f\xrightarrow{\kappa}M
\xrightarrow{q}I\longrightarrow0,
\qquad
0\longrightarrow I\xrightarrow{j}N
\xrightarrow{p}\operatorname{coker}f\longrightarrow0
\]
使用正合性。第二个序列说明 \(F(j)\) 单射，因而
\[
\ker F(f)=\ker\bigl(F(j)F(q)\bigr)=\ker F(q).
\]
第一个序列说明 \(F(\kappa)\) 单射，且其像恰为 \(\ker F(q)\)。因此 \(F(\kappa):F(\ker f)\rightarrow F(M)\) 是 \(F(f)\) 的核映射，\(F\) 保持 \(f\) 的核。

另一方面，第一个序列说明 \(F(q)\) 满射，因此
\[
\operatorname{im}F(f)=\operatorname{im}\bigl(F(j)F(q)\bigr)
=\operatorname{im}F(j).
\]
第二个序列说明 \(F(p)\) 满射，且其核恰为 \(\operatorname{im}F(j)\)，从而给出
\[
F(N)/\operatorname{im}F(j)\cong F(\operatorname{coker}f),
\]
这个同构由 \(F(p)\) 诱导，所以 \(F(p)\) 正是 \(F(f)\) 的余核映射，\(F\) 也保持 \(f\) 的余核。又因 \(F(f-g)=F(f)-F(g)\)，它保持等化子与余等化子。最后应用\cref{b4:thm:limits-products-equalizers}及其对偶构造，得到所述有限极限与有限余极限的保持性。
\end{proof}

\begin{example}[正合函子不保持无限积]\label{b4:exm:exact-not-products}
函子 \(F:\mathbf{Ab}\rightarrow\mathbb Q\text{-}\mathbf{Mod}\)，
\(F(M)=\mathbb Q\otimes_{\mathbb Z}M\)，是正合的，却不保持可数积。先说明正合性所用的事实：这个张量积可写成分母为非零整数的分式模，且 \(m/s=0\) 当且仅当某个非零整数零化 \(m\)。张量积到分式模的映射为 \((a/b)\otimes m\mapsto am/b\)，逆映射为 \(m/s\mapsto(1/s)\otimes m\)，由平衡关系可验证互逆。对于短正合列 \(0\rightarrow A\rightarrow B\rightarrow C\rightarrow0\)，分式模中的单射性由上述零元判据得到；若 \(b/s\) 在 \(C\) 中为零，则某个非零整数 \(t\) 使 \(tb\) 来自 \(A\)，从而 \(b/s=(tb)/(ts)\) 来自 \(A\)；满射性则逐个提升分子。这证明了 \(F\) 正合。

\ProseParagraphBreak
现在考虑比较映射
\[
\mathbb Q\otimes_{\mathbb Z}\prod_{n\ge1}\mathbb Z
\longrightarrow\prod_{n\ge1}\mathbb Q.
\]
左边每个元素都可通分为一个 \((a_n)/s\)，故其像具有统一的非零整数分母。元组 \((2^{-n})_{n\ge1}\) 没有这种分母，因此不在像中。失败的是无限积；有限个分母总能通分，与上一命题并不矛盾。
\end{example}

\begin{example}[正合函子不保持无限余积]\label{b4:exm:exact-not-coproducts}
设 \(k\) 为域，取 \(F:k\text{-}\mathbf{Mod}\rightarrow k\text{-}\mathbf{Mod}\)，\(F(V)=\prod_{n\ge1}V\)，态射逐分量作用。这个函子是加性的，并且正合：向量空间的每个短正合列都分裂，选定的线性截面逐分量作用，就给出施加 \(F\) 后的截面；核仍逐分量计算。

\ProseParagraphBreak
对于可数族 \(V_j=k\)，余积的比较映射为
\[
\bigoplus_{j\ge1}\prod_{n\ge1}k
\longrightarrow
\prod_{n\ge1}\bigoplus_{j\ge1}k.
\]
左边元素只有有限多个非零的 \(j\) 分量，所以像中的所有行在同一个有限的列集合内非零。右边允许每一行分别具有有限支集。矩阵 \(a_{n,j}=\delta_{n,j}\) 属于右边，却不在左边的像中，故比较映射不是满射。
\end{example}

由此，模范畴中的正合性控制核、余核与有限直和，却不能自动控制无限积或无限余积。若一个正合函子保持所有小积，则它保持所有小极限；若它保持所有小余积，则它保持所有小余极限。两者都由积与等化子、余积与余等化子的构造判据得到。


\begin{chaptersummary}
极限刻画从同一对象出发、与图表相容的态射族；余极限则刻画从图表各对象通向同一对象的相容态射族。由积与等化子构造极限、由余积与余等化子构造余极限，把任意小图表的泛性质归结为两种基本构造。在集合、群、交换环和模中，余极限常通过先取余积、再施加相容关系来实现。在模中，两种构造分别成为直积之间一个同态的核与直和之间一个同态的余核。

\ProseParagraphBreak
滤过性保证有限多个元素和等式能够在一个阶段内处理。这既给出滤过余极限的正合性，也给出它与有限极限的交换。无限多个条件却可能没有共同阶段，因此不能据此推出与无限积交换。在模范畴之间，加性左正合函子保持有限极限，加性右正合函子保持有限余极限，所以加性正合函子同时保持两者；有理化函子及可数积函子分别说明，正合性仍不足以保证保持无限积或无限余积。
\end{chaptersummary}

\BookExercises
以下环均含幺，模均为含幺左模，环同态保持单位元。未特别说明时，指标范畴为小范畴。

\begin{exercise}\label{b4:ex:02-equalizer-mono}
设范畴中的态射 \(i:E\rightarrow X\)、\(r:X\rightarrow E\) 满足 \(ri=1_E\)。证明 \(i\) 是 \(ir,1_X:X\rightrightarrows X\) 的等化子。对偶地，若 \(p:Y\rightarrow Q\)、\(s:Q\rightarrow Y\) 满足 \(ps=1_Q\)，证明 \(p\) 是 \(sp,1_Y:Y\rightrightarrows Y\) 的余等化子。最后，在交换环范畴中证明包含映射 \(\mathbb Z\rightarrow\mathbb Q\) 是单态射，却不可能为任何一对态射的等化子。
\end{exercise}
\begin{solution}
由 \(iri=i\)，态射 \(i\) 确实使 \(ir,1_X\) 相等。若 \(h:T\rightarrow X\) 满足 \(irh=h\)，则 \(rh:T\rightarrow E\) 是所需分解；任何分解 \(h=ik\) 都满足 \(k=rik=rh\)，故唯一。对偶地，\(psp=p\)。若 \(h:Y\rightarrow T\) 满足 \(hsp=h\)，则唯一分解为 \(h=(hs)p\)，唯一性由后复合 \(s\) 得到。

整数到有理数的包含是底集合单射，因而在交换环范畴中是单态射。对任意交换环 \(A\)，至多存在一个保幺同态 \(\mathbb Q\rightarrow A\)：整数的像由单位元决定，非零整数的逆若存在则唯一，有理数的像也随之确定。因此任意平行态射 \(f,g:\mathbb Q\rightarrow A\) 必相等，它们的等化子是 \(1_{\mathbb Q}\)。等化子的唯一性将迫使任何其他等化子到 \(\mathbb Q\) 的结构映射为同构，而 \(\mathbb Z\rightarrow\mathbb Q\) 不是同构，故不可能是等化子。这说明单态射未必能由等化条件描述。
\end{solution}

\begin{exercise}\label{b4:ex:02-pullback-intersection}
设 \(N,P\) 是 \(R\) 模 \(M\) 的子模。证明包含映射 \(N\rightarrow M\leftarrow P\) 的纤维积同构于 \(N\cap P\)，而商映射 \(M/N\leftarrow M\rightarrow M/P\) 的纤维余积同构于 \(M/(N+P)\)。
\end{exercise}
\begin{solution}
纤维积中的元素是满足 \(n=p\) 的 \((n,p)\)，映射 \(x\mapsto(x,x)\) 给出 \(N\cap P\) 到它的同构，并保持投影。对于纤维余积，两个商模到 \(T\) 的同态若在 \(M\) 上一致，等价于一个 \(M\rightarrow T\) 的同态同时零化 \(N,P\)，即零化 \(N+P\)。这又等价于唯一的 \(M/(N+P)\rightarrow T\)，故得到所需泛性质。
\end{solution}

\begin{exercise}\label{b4:ex:02-integer-pushout}
设 \(m,n\) 为正整数。计算同态 \(\mathbb Z\xrightarrow{\times m}\mathbb Z\)、\(\mathbb Z\xrightarrow{\times n}\mathbb Z\) 的模纤维余积，并求两个结构映射的核。
\end{exercise}
\begin{solution}
纤维余积为 \(Q=\mathbb Z^2/\mathbb Z(m,-n)\)。令 \(d=\gcd(m,n)\)，取整数 \(a,b\) 使 \(am+bn=d\)。矩阵
\[
U=\begin{pmatrix}a&-b\\ n/d&m/d\end{pmatrix}
\]
的行列式为一，并把 \((m,-n)^{\mathsf T}\) 送到 \((d,0)^{\mathsf T}\)，故 \(Q\cong\mathbb Z/d\mathbb Z\oplus\mathbb Z\)。两个结构映射的核都为零：若 \((x,0)=t(m,-n)\)，则 \(tn=0\)，所以 \(t=x=0\)；另一映射同理。这里结构映射单射是该具体计算的结果，不能作为一般纤维余积的性质。
\end{solution}

\begin{exercise}\label{b4:ex:02-ring-coproduct}
在交换环范畴中计算 \(\mathbb Z/m\mathbb Z\) 与 \(\mathbb Z/n\mathbb Z\) 的余积，其中 \(m,n\) 为正整数。说明它在什么条件下为零环。
\end{exercise}
\begin{solution}
一个同时接收这两个环的保幺映射的交换环，必须满足 \(m\cdot1=n\cdot1=0\)，这等价于 \(d\cdot1=0\)，其中 \(d=\gcd(m,n)\)。因此这些相容映射唯一来自 \(\mathbb Z/d\mathbb Z\) 的保幺映射，余积就是 \(\mathbb Z/d\mathbb Z\)，也即相应的整数张量积。它为零环当且仅当 \(d=1\)。
\end{solution}

\begin{exercise}\label{b4:ex:02-finite-completeness}
证明一个范畴若有终对象和所有纤维积，就有所有有限极限。写出等化子的纤维积构造。再陈述并证明对偶结论。
\end{exercise}
\begin{solution}
设终对象为 \(1\)。\(X\rightarrow1\leftarrow Y\) 的纤维积就是 \(X,Y\) 的积，迭代可得任意有限积，包括空积。对 \(f,g:X\rightarrow Y\)，取
\[
X\xrightarrow{(f,g)}Y\times Y
\ \xleftarrow{\Delta}\ Y,
\]
其中 \(\Delta\) 的两个分量均为 \(1_Y\)。其纤维积到 \(X\) 的结构映射是等化子：相容条件恰为 \(fh=gh\)，而到 \(Y\) 的另一映射随后由这个共同值唯一决定。对有限指标范畴，\cref{b4:thm:limits-products-equalizers}中的两个积都只含有限个因子，因此所有有限极限存在。

反转箭头可得：有始对象和所有纤维余积的范畴具有所有有限余极限。具体地，先以始对象为公共来源得到二元余积；对 \(f,g:X\rightarrow Y\)，沿余对角映射 \(X\amalg X\rightarrow X\) 与在两个余积项上分别为 \(f,g\) 的映射 \(X\amalg X\rightarrow Y\) 取纤维余积，就得到余等化子。
\end{solution}

\begin{exercise}\label{b4:ex:02-initial-index}
设指标范畴 \(J\) 有始对象 \(0\)，\(D:J\rightarrow\mathcal C\) 为图表。证明 \(\lim_JD\cong D(0)\)。若 \(J\) 有终对象 \(1\)，证明 \(\operatorname{colim}_JD\cong D(1)\)。解释为何不可随意交换“始”与“终”。
\end{exercise}
\begin{solution}
记唯一的态射 \(0\rightarrow j\) 为 \(t_j\)。对任意 \(\alpha:j\rightarrow k\)，始对象的唯一性给出 \(\alpha t_j=t_k\)，所以 \(D(\alpha)D(t_j)=D(t_k)\)。这正说明 \((D(t_j):D(0)\rightarrow D(j))_j\) 是锥。若另有锥 \((c_j:T\rightarrow D(j))_j\)，相容性给出 \(c_j=D(t_j)c_0\)，故 \(c_0\) 是到上述锥的分解；因 \(t_0=1_0\)，任何分解的 \(0\) 分量又迫使它等于 \(c_0\)，所以分解唯一。第二个结论对偶。若 \(J\) 只有一条非恒等箭头 \(0\rightarrow1\)，则图表就是 \(X\rightarrow Y\)：其极限为 \(X\)，余极限为 \(Y\)。取 \(X,Y\) 不同构即可看出两种说法不能交换。
\end{solution}

\begin{exercise}\label{b4:ex:02-group-action-diagram}
设 \(G\) 为群，把它看成只有一个对象、态射为 \(G\) 的范畴 \(BG\)。对一个左 \(G\) 集 \(X\)，将其视为图表 \(BG\rightarrow\mathbf{Set}\)。计算该图表的极限与余极限。
\end{exercise}
\begin{solution}
锥只有一个映射 \(T\rightarrow X\)，相容条件为其像中每个点都被所有 \(g\in G\) 固定。因此极限是全体不动点 \(X^G\)。一个余锥是满足 \(f(gx)=f(x)\) 的映射 \(X\rightarrow T\)，即在每条轨道上常值的映射；它唯一经过轨道集 \(G\backslash X\) 分解，所以余极限为轨道集。
\end{solution}

\begin{exercise}\label{b4:ex:02-parallel-filtered-condition}
令 \(J\) 只有两个对象 \(0,1\)，除恒等态射外恰有两条不同的箭头 \(a,b:0\rightrightarrows1\)。构造集合图表，说明缺少平行态射条件时，余极限中相等的两个同阶段元素未必能在某个后续对象中变成相等。再向 \(J\) 添加对象 \(2\) 和箭头 \(c:1\rightarrow2\)，并规定 \(ca=cb\)，所得范畴 \(J'\) 除这些箭头、共同复合 \(ca\) 及恒等态射外没有其他态射。证明 \(J'\) 滤过，并将所构造的图表延拓到 \(J'\)，检验新增的过渡映射怎样使上述两个元素变成相等。
\end{exercise}
\begin{solution}
\(J\) 正是\cref{b4:def:filtered-category}之后的平行箭头例子：任意两个对象有共同目标，但两条平行箭头不能通过后复合变成相等。取 \(D(0)=\{*\}\)、\(D(1)=\{x,y\}\)，令 \(D(a)(*)=x,D(b)(*)=y\)。余极限把 \(x,y\) 识别，所以是单点集；但从对象 \(1\) 出发只有恒等态射，没有能使 \(x,y\) 相等的过渡映射。

在 \(J'\) 中，对象 \(2\) 是所有对象的共同目标；唯一一对不同的平行态射仍为 \(a,b\)，而 \(c\) 使它们的后复合相等，所以 \(J'\) 滤过。令 \(D'(2)=\{z\}\)，令 \(D'(c)\) 将 \(x,y\) 都送到 \(z\)，并在 \(J\) 上保留原图表。于是 \(D'(c)D'(a)=D'(c)D'(b)\)，确实得到延拓。新余极限仍为单点集，但这次 \(x,y\) 的相等可由同一条过渡映射 \(D'(c)\) 直接见证。共同目标只容许把元素送到同一对象；平行态射条件还保证比较这些元素的不同路径能够统一。
\end{solution}

\begin{exercise}\label{b4:ex:02-directed-quotients}
设 \((N_i)_{i\in I}\) 是 \(R\) 模 \(M\) 的有向递增子模族。证明
\[
\varinjlim_i M/N_i\cong M/\bigcup_iN_i.
\]
若正向系统 \(M_0\rightarrow M_1\rightarrow\cdots\) 的所有过渡映射 \(M_n\rightarrow M_{n+1}\) 均为零，计算其正向极限。
\end{exercise}
\begin{solution}
令 \(N=\bigcup_iN_i\)。从每个商 \(M/N_i\) 到 \(M/N\) 的映射相容且满射。若 \(m+N_i\) 被送到零，则 \(m\in N_j\) 对某个 \(j\) 成立，取 \(i,j\) 的共同上界 \(k\)，它在 \(M/N_k\) 中已经为零。因此诱导映射为单射，也为满射。第二问中，每个阶段的元素映到下一阶段即为零，所以元素判据给出余极限为零。
\end{solution}

\begin{exercise}\label{b4:ex:02-inverse-multiplication}
计算逆向系统
\[
\cdots\xrightarrow{\times2}\mathbb Z
\xrightarrow{\times2}\mathbb Z
\xrightarrow{\times2}\mathbb Z
\]
的逆极限。再证明，当整数模 \(M\) 满足乘二映射为同构时，同样的逆向系统的逆极限同构于 \(M\)。
\end{exercise}
\begin{solution}
相容元组满足 \(a_n=2a_{n+1}\)。故每个 \(a_n\) 被任意高次的二整除，只能为零，逆极限为零。若 \(2:M\rightarrow M\) 可逆，则任选 \(m_0\in M\)，其后各项唯一为 \(m_n=2^{-n}m_0\)。首项投影给出所求同构。两种情况说明，过渡映射与对象同样决定逆极限。
\end{solution}

\begin{exercise}\label{b4:ex:02-modular-projections}
设 \(p\) 为素数。证明每个投影 \(\mathbb Z_p\rightarrow\mathbb Z/p^n\mathbb Z\) 都是满射，且其核为 \(p^n\mathbb Z_p\)。
\end{exercise}
\begin{solution}
任意模 \(p^n\) 的类可取整数代表，这个整数的全部剩余类构成一个提升，所以投影满射。设 \(x\in\mathbb Z_p\) 的数字为 \(c_0,c_1,\ldots\)，其中 \(0\le c_i<p\)。它落在核中，等价于 \(c_0=\cdots=c_{n-1}=0\)。令 \(y\in\mathbb Z_p\) 的数字为 \(c_n,c_{n+1},\ldots\)。若 \(m\le n\)，则 \(x\) 与 \(p^ny\) 在模 \(p^m\) 的商中均为零；若 \(m>n\)，则它们在该商中的代表满足
\[
\sum_{i=n}^{m-1}c_ip^i
=p^n\sum_{j=0}^{m-n-1}c_{n+j}p^j.
\]
所以 \(x=p^ny\) 在每个有限商中成立，进而在逆极限中成立。反过来，任何 \(p^ny\) 在模 \(p^n\) 的商中都为零，故投影的核恰为 \(p^n\mathbb Z_p\)。
\end{solution}

\begin{exercise}\label{b4:ex:02-filtered-kernels}
设 \(J\) 为滤过范畴，\(f:M\Rightarrow N\) 为 \(R\) 模值图表之间的自然变换。由本章的交换定理说明
\[
\operatorname{colim}_{j\in J}\ker f_j
\cong
\ker\bigl(\operatorname{colim}_JM\rightarrow
\operatorname{colim}_JN\bigr).
\]
成立。再取 \(J=BC_2\)，令 \(M,N\) 都为整数模，\(C_2\) 的非单位元均作用为乘负一，\(f:M\Rightarrow N\) 为乘二。计算等式两边，证明去掉滤过性后结论可能失败。
\end{exercise}
\begin{solution}
核是一个态射与零态射的等化子，所以滤过余极限与有限极限的交换给出第一式。

对反例，乘二与乘负一交换，故 \(f\) 确实是自然变换，其逐项核为零，等式左边为零。一个模值 \(BC_2\) 图表的余极限，是把每个 \(m\) 与其作用后的像识别所得的商模。本题作用为乘负一，故两份余极限均为 \(\mathbb Z/2\mathbb Z\)。乘二诱导其上的零映射，因而等式右边的核为整个 \(\mathbb Z/2\mathbb Z\)，不等于左边。\(BC_2\) 不是滤过范畴：单位态射与非单位态射不能通过后复合变为相等，因为它的每个态射均可逆。
\end{solution}

\begin{exercise}\label{b4:ex:02-hom-finite-presentation}
设 \(P\) 是有限展示的左 \(R\) 模，\(M:J\rightarrow R\text{-}\mathbf{Mod}\) 是滤过图表。证明自然映射
\[
\operatorname{colim}_{j\in J}\operatorname{Hom}_R(P,M_j)
\longrightarrow
\operatorname{Hom}_R(P,\operatorname{colim}_JM)
\]
为双射。指出有限生成元与有限关系分别用在何处。
\end{exercise}
\begin{solution}
选有限展示 \(R^a\rightarrow R^b\rightarrow P\rightarrow0\)。一个 \(P\rightarrow M\) 的同态等价于 \(b\) 个满足由 \(a\) 个关系给出的有限组线性等式的元素。取值在余极限中的同态，其有限个生成元的像可同时提升到某个 \(M_i\)；有限个关系在余极限中成立，故再送到某个后续 \(M_j\) 后同时成立，得到同态 \(P\rightarrow M_j\)，证明满射性。若两个阶段的同态在余极限中相等，先送到共同阶段，再使有限个生成元的像同时相等，两个同态在那里就相等，证明单射性。共同提升像和检验同态相等只需有限生成元；使全部定义关系同时成立还需要有限关系。
\end{solution}

\begin{exercise}\label{b4:ex:02-filtered-empty}
说明滤过范畴的非空要求不能从“滤过余极限与有限极限交换”的集合版本中删除。另说明：有限极限交换命题只给出充分条件，并不声称指标范畴必须只有有限条态射。
\end{exercise}
\begin{solution}
取余极限指标为空范畴、极限指标也为空范畴。先取空极限得到单点集，再对空图表取余极限得到空集；先取余极限后取空极限，则得到单点集。比较映射 \(\varnothing\rightarrow\{*\}\) 不为同构。

对于第二句，取有限生成群 \(G=\langle g_1,\ldots,g_r\rangle\)，把它视为 \(BG\)。对任意 \(G\) 集 \(X\)，有
\[
X^G=\{x\in X:g_ix=x\ (1\le i\le r)\}
=\operatorname{Eq}\!\left(
X\underset{x\mapsto(x,\ldots,x)}{\overset{x\mapsto(g_1x,\ldots,g_rx)}{\rightrightarrows}}X^r
\right).
\]
这是因为固定每个生成元的点也固定它们的逆及任意乘积。这个构造只使用有限积和等化子，因而与滤过余极限交换。\(r=0\) 时 \(G\) 为平凡群，\(X^0\) 为单点集，两条映射自动相等，仍得到 \(X^G=X\)。所以即使 \(G\) 无限、\(BG\) 有无限条态射，\(BG\) 形极限也可以与滤过余极限交换。
\end{solution}

\begin{exercise}\label{b4:ex:02-infinite-fixed-points}
令 \(G=\bigoplus_{r\ge1}C_2\)，\(H_n=\bigoplus_{1\le r\le n}C_2\)，其中 \(n\ge1\)。把 \(G\) 视为只有一个对象 \(*\) 的范畴 \(BG\)，并以正整数的通常次序定义指标范畴 \(\mathbb N_{>0}\)。令
\[
D:\mathbb N_{>0}\times BG\longrightarrow\mathbf{Set},
\qquad D(n,*)=G/H_n,
\]
其中群元素作用为左平移，\(n\le m\) 时的过渡映射为商映射 \(G/H_n\rightarrow G/H_m\)。证明这些映射组成一个图表，其正向余极限为单点集，但每个 \(G/H_n\) 都没有 \(G\) 不动点。由此写出滤过余极限不与 \(BG\) 形极限交换的例子。
\end{exercise}
\begin{solution}
商映射保持左平移，故两个方向的态射交换；各方向又都保持恒等与复合，因此确实定义了图表 \(D\)。任意两个陪集的代表之差属于某个 \(H_N\)，所以它们最终在同一商中相等，余极限为单点集。若一个陪集被全部 \(G\) 固定，则其稳定子群应为 \(G\)，但因 \(G\) 交换，该稳定子群实际上为真子群 \(H_n\)，故没有不动点。由\cref{b4:ex:02-group-action-diagram}，不动点就是 \(BG\) 形极限。因此
\[
\begin{aligned}
\varinjlim_n\bigl(\lim_{BG}D(n,-)\bigr)&=\varnothing,\\
\lim_{BG}\bigl(\varinjlim_nD(n,-)\bigr)&=\{*\}.
\end{aligned}
\]
对应的比较映射 \(\varnothing\rightarrow\{*\}\) 不是同构。
\end{solution}

\begin{exercise}\label{b4:ex:02-hom-coproduct}
设 \(R\) 为含幺环，\((M_i)_{i\in I}\) 为左模族，\(N\) 为左模。证明
\[
\operatorname{Hom}_R\!\left(\bigoplus_iM_i,N\right)
\cong\prod_i\operatorname{Hom}_R(M_i,N).
\]
当所有模都是域 \(k\) 上的向量空间时，以 \(M_i=k\) 为例，说明右边一般不能改成直和。
\end{exercise}
\begin{solution}
映射送到它在每个余积项上的限制；给定右边的任意同态族 \((f_i)\)，逆映射为 \(\sum_i m_i\mapsto\sum_i f_i(m_i)\)，因为输入有有限支集，所以输出总有意义。若 \(I\) 无限、\(M_i=N=k\)，在每个分量上取恒等映射得到合法的同态族，却不具有有限支集。因此它属于右边的积而不属于直和。
\end{solution}

\begin{exercise}\label{b4:ex:02-rationalization-product}
精确描述自然映射
\[
\mathbb Q\otimes_{\mathbb Z}\prod_{n\ge1}\mathbb Z
\longrightarrow\prod_{n\ge1}\mathbb Q
\]
的像，并证明该映射是单射。
\end{exercise}
\begin{solution}
张量积中的每个元素都是有限和 \(\sum_{i=1}^r(r_i/s_i)\otimes(a_n^{(i)})\)，其中 \(r_i,s_i\) 为整数、\(s_i\ne0\)。取这些分母的非零公倍数 \(s\)，利用整数张量积的平衡关系，便可通分合并为
\[
\sum_{i=1}^r\dfrac{r_i}{s_i}\otimes(a_n^{(i)})
=\dfrac1s\otimes\left(\sum_{i=1}^r r_i\dfrac{s}{s_i}a_n^{(i)}\right)_n.
\]
把这种元素简记为 \((a_n)/s\)，它的像就是 \((a_n/s)_n\)。因此像恰为具有统一非零整数分母的有理数族，即存在非零整数 \(s\)，使所有 \(sq_n\) 都为整数的族。反向包含取 \((1/s)\otimes(sq_n)_n\) 即得。若 \((a_n)/s\) 映到零，则每个 \(a_n/s=0\) 于 \(\mathbb Q\)，从而全部 \(a_n=0\)，原元素为 \((1/s)\otimes0=0\)。这证明单射。因而此反例的障碍完全是满射性。
\end{solution}

\begin{exercise}\label{b4:ex:02-exact-product-functor}
设 \(k\) 为域，\(F(V)=V^{\mathbb N_{>0}}\)。证明 \(F\) 保持所有小积，却不保持可数余积。后一个结论中，精确描述余积比较映射的像。
\end{exercise}
\begin{solution}
对任意集合指标族 \((V_j)\)，两种积都有同一个分量集合：
\[
F\!\left(\prod_jV_j\right)
\cong\prod_{(n,j)}V_j
\cong\prod_jF(V_j).
\]
这些对应保持投影，故就是积的比较同构。取全部 \(V_j=k\)，余积比较映射的像由满足以下条件的矩阵组成：存在同一个有限集合 \(S\subseteq\mathbb N_{>0}\)，使所有 \(j\notin S\) 的列均为零。目标仅要求每一行分别只有有限个非零项。对角矩阵不满足统一有限列条件，故映射不满射。
\end{solution}

\begin{exercise}\label{b4:ex:02-limits-exact-criterion}
设 \(R,S\) 为含幺环，\(F:R\text{-}\mathbf{Mod}\rightarrow S\text{-}\mathbf{Mod}\) 是幺性左模范畴之间的加性函子。证明：若 \(F\) 左正合且保持所有小积，则保持所有小极限；若 \(F\) 右正合且保持所有小余积，则保持所有小余极限。
\end{exercise}
\begin{solution}
由\cref{b4:prop:exact-preserves-finite}，加性左正合函子保持有限极限，特别保持等化子。再配合所有小积，\cref{b4:thm:limits-products-equalizers}把每个小图表的极限写成两个小积之间的等化子，故它保持所有小极限。对偶地，加性右正合函子保持余等化子；配合所有小余积，再用该定理的对偶构造即可。这里增加的积或余积保持条件，使\cref{b4:prop:exact-preserves-finite}的有限结论推广到任意小图表。零对象由加性保证，空积与空余积也包括在内。
\end{solution}

\begin{exercise}\label{b4:ex:02-inverse-left-exact}
设 \(R\) 为含幺环，以正整数为指标的幺性左 \(R\) 模逆向系统在每个阶段组成短正合列
\[
0\longrightarrow A_n\xrightarrow{u_n}B_n
\xrightarrow{v_n}C_n\longrightarrow0,
\]
其中过渡同态 \(\rho_n^A:A_{n+1}\rightarrow A_n\)、\(\rho_n^B:B_{n+1}\rightarrow B_n\)、\(\rho_n^C:C_{n+1}\rightarrow C_n\) 与 \(u_n,v_n\) 相容。若每个 \(\rho_n^A\) 都是满射，证明
\[
0\longrightarrow\varprojlim_nA_n
\longrightarrow\varprojlim_nB_n
\longrightarrow\varprojlim_nC_n\longrightarrow0
\]
也是短正合列。说明如何逐次修正 \(C_n\) 中元素的提升，使它们相容。
\end{exercise}
\begin{solution}
前两处正合由\cref{b4:prop:inverse-limits-left-exact}得到。给定 \((c_n)\in\varprojlim_nC_n\)，先选 \(b_1\in B_1\)，使 \(v_1(b_1)=c_1\)。假设已经选定 \(b_n\)，在下一阶段任取 \(\widetilde b_{n+1}\in B_{n+1}\)，使 \(v_{n+1}(\widetilde b_{n+1})=c_{n+1}\)。我们要修正这个任意提升，使它过渡到 \(b_n\)，又不能改变它在 \(C_{n+1}\) 中的像。因此修正项必须属于 \(\ker v_{n+1}=\operatorname{im}u_{n+1}\)。另一方面，下一阶段提升过渡后与 \(b_n\) 的差属于 \(\ker v_n=\operatorname{im}u_n\)；\(\rho_n^A\) 的满射性恰好允许把这个差从 \(A_n\) 提升到 \(A_{n+1}\)。具体地，由相容性，
\[
v_n\bigl(b_n-\rho_n^B(\widetilde b_{n+1})\bigr)
=c_n-\rho_n^C(c_{n+1})=0,
\]
故存在 \(a_n\in A_n\)，使 \(u_n(a_n)=b_n-\rho_n^B(\widetilde b_{n+1})\)。由 \(\rho_n^A\) 满射，可取 \(a_{n+1}'\in A_{n+1}\) 满足 \(\rho_n^A(a_{n+1}')=a_n\)。令
\[
b_{n+1}=\widetilde b_{n+1}+u_{n+1}(a_{n+1}').
\]
此时 \(v_{n+1}(b_{n+1})=c_{n+1}\)，且 \(\rho_n^B(b_{n+1})=b_n\)。递归得到相容族 \((b_n)\)，它映到给定的 \((c_n)\)，所以逆极限之间的最后一个映射也是满射。
\end{solution}
